% ===============================================================
%  Characterizing Speculative Decoding Dynamics for LLMs
%  on Consumer-Class GPUs  --  ICECCME 2026 conference talk
%
%  Target duration: 8-12 minutes (planned ~10:40, see \note pages)
%
%  BUILD
%    Slides only :  latexmk -pdf slides.tex
%    Notes only  :  latexmk -pdf -jobname=slides-notes \
%                     "\def\NOTESMODE{}% ===============================================================
%  Characterizing Speculative Decoding Dynamics for LLMs
%  on Consumer-Class GPUs  --  ICECCME 2026 conference talk
%
%  Target duration: 8-12 minutes (planned ~10:40, see \note pages)
%
%  BUILD
%    Slides only :  latexmk -pdf slides.tex
%    Notes only  :  latexmk -pdf -jobname=slides-notes \
%                     "\def\NOTESMODE{}% ===============================================================
%  Characterizing Speculative Decoding Dynamics for LLMs
%  on Consumer-Class GPUs  --  ICECCME 2026 conference talk
%
%  Target duration: 8-12 minutes (planned ~10:40, see \note pages)
%
%  BUILD
%    Slides only :  latexmk -pdf slides.tex
%    Notes only  :  latexmk -pdf -jobname=slides-notes \
%                     "\def\NOTESMODE{}% ===============================================================
%  Characterizing Speculative Decoding Dynamics for LLMs
%  on Consumer-Class GPUs  --  ICECCME 2026 conference talk
%
%  Target duration: 8-12 minutes (planned ~10:40, see \note pages)
%
%  BUILD
%    Slides only :  latexmk -pdf slides.tex
%    Notes only  :  latexmk -pdf -jobname=slides-notes \
%                     "\def\NOTESMODE{}\input{slides.tex}"
%    Dual screen :  set \def\NOTESMODE{} -> change option below
% ===============================================================
\documentclass[aspectratio=169, 11pt]{beamer}

\input{preamble}

% --- notes switch -----------------------------------------------
\ifdefined\NOTESMODE
  \setbeameroption{show only notes}
  % Notes-only build: drop the slide thumbnail and shrink the type so that
  % even the longest note fits on a single page.
  \setbeamertemplate{note page}[plain]
  \setbeamerfont{note page}{size=\footnotesize}
\else
  \setbeameroption{hide notes}
  % For a presenter-display setup, comment the line above and use:
  % \setbeameroption{show notes on second screen=right}
\fi
% ----------------------------------------------------------------

\title{Characterizing Speculative Decoding Dynamics for Large
       Language Models on Consumer-Class GPUs}
\author{Samsu Sempena \and Zhuohan Cui}
\institute{The University of Sydney}
\date{ICECCME 2026 \\ 15--17 October 2026, Bali, Indonesia}

\begin{document}

% ================================================================
\begin{frame}[plain]
  \titlepage
\end{frame}
\note{%
\textbf{[0:00--0:20 | Title]}\\[2pt]
``Good morning. I'm Samsu Sempena from the University of Sydney, and this is
joint work with Zhuohan Cui. Our paper characterizes what speculative decoding
actually does on consumer-class GPUs --- the kind of hardware a practitioner
has under their desk, not in a data centre.''\\[4pt]
\emph{Cue:} say the one-line thesis up front so the audience has a hook ---
``the short version is that the \emph{policy ranking} travels between GPUs,
but the \emph{speedup} does not.''\\[4pt]
\emph{Housekeeping:} confirm the clicker works, and check whether the chair
wants questions during or after.
}

% ================================================================
\section{Motivation}

\begin{frame}{Decoding is sequential, and that is the bottleneck}
  \begin{columns}[T,onlytextwidth]
    \begin{column}{0.56\textwidth}
      \begin{itemize}
        \item Autoregressive decoding needs \alert{one full forward pass of the
              target model per token}.
        \item Single-stream, interactive latency is therefore bound by
              \emph{memory traffic}, not by arithmetic.
        \item Speculative decoding is the standard fix --- and it is reported
              to be a large, free win.
      \end{itemize}
      \vspace{4pt}
      \begin{block}{The practitioner's problem}
        Almost all published gains are measured on \textbf{data-centre}
        hardware. On a single RTX-class GPU, what actually happens?
      \end{block}
    \end{column}
    \begin{column}{0.42\textwidth}
      \scriptsize
      \begin{tabular}{@{}ll@{}}
        \toprule
        \multicolumn{2}{l}{\textbf{Open deployment questions}}\\
        \midrule
        Which draft size?      & 0.5B? 1.5B? \\
        Which proposal length? & $k=4$? $8$? $16$? \\
        Greedy or sampling?    & speed vs.\ quality \\
        Does it transfer?      & desktop $\to$ laptop \\
        \bottomrule
      \end{tabular}
    \end{column}
  \end{columns}
\end{frame}
\note{%
\textbf{[0:20--1:05 | Motivation]}\\[2pt]
Frame the bottleneck physically: every token costs a complete pass over the
target model's weights, so on a single stream you are moving gigabytes of
weights per token and the GPU's compute units are mostly idle. That is why
batch-1 interactive decoding is memory-bandwidth bound.\\[4pt]
Then the pivot: ``speculative decoding is the well-known answer, and the
literature reports two-to-five-times speedups. But look at where those numbers
come from --- A100s, H100s, B200s. The practitioner running a 3B or 4B model on
one RTX card has no reliable guidance.''\\[4pt]
Point at the four questions on the right --- these are exactly the knobs
someone has to set on day one, and say we will answer all four with
measurements.\\[4pt]
\emph{Do not} over-explain memory bandwidth; one sentence is enough.
}

% ================================================================
\begin{frame}{Speculative decoding in one slide}
  \begin{columns}[T,onlytextwidth]
    \begin{column}{0.50\textwidth}
      \textbf{Draft $\to$ verify $\to$ commit}
      \begin{enumerate}\itemsep2pt
        \item A small \textbf{draft} model proposes $k$ tokens
              autoregressively.
        \item The large \textbf{target} verifies all $k$ in
              \alert{one} parallel forward pass.
        \item Commit the accepted prefix; add one correction token on the
              first mismatch.
      \end{enumerate}
      \vspace{3pt}
      Modified rejection sampling keeps the output distribution
      \emph{identical} to the target's.
      \\[2pt]
      \tiny Leviathan et al.\ 2023; Chen et al.\ 2023
    \end{column}
    \begin{column}{0.48\textwidth}
      \begin{block}{Per-cycle latency}
        \[
          L \;=\; \frac{T_{\text{draft}} + T_{\text{verify}}}{\tau},
          \qquad
          T_{\text{draft}} = k \cdot t_{\text{draft step}}
        \]
        $\tau \in [1,\,k{+}1]$ accepted tokens per cycle.
      \end{block}
      \vspace{-2pt}
      \begin{itemize}\itemsep2pt
        \item Numerator grows \alert{linearly} in $k$.
        \item Denominator $\tau$ saturates as acceptance $\alpha$ falls.
        \item[$\Rightarrow$] There must be a $k$ beyond which it \emph{loses}.
      \end{itemize}
    \end{column}
  \end{columns}
\end{frame}
\note{%
\textbf{[1:05--2:00 | Method recap]}\\[2pt]
Keep this brisk --- most of the room knows the algorithm. Walk the three steps
once, then land hard on the equation, because it is the analytical spine of
every result that follows.\\[4pt]
Say it plainly: ``the cost in the numerator grows linearly with $k$, because
you run the draft model $k$ times. The payoff in the denominator is $\tau$, the
accepted tokens per cycle, and $\tau$ saturates --- once acceptance drops off,
extra proposals are wasted work. So the theory already tells you there is a
crossover. Our contribution is showing \emph{where} it sits on real consumer
hardware, and that it sits much lower than people assume.''\\[4pt]
Stress the correctness guarantee once --- ``this is lossless with respect to
the target distribution'' --- because it pre-empts a question, and because it
makes the quality result later more interesting.
}

% ================================================================
\begin{frame}{What is missing, and what we ask}
  \begin{columns}[T,onlytextwidth]
    \begin{column}{0.44\textwidth}
      \textbf{Three gaps in the literature}
      \begin{itemize}\itemsep3pt
        \item No cross-model \emph{and} cross-hardware validation at the
              consumer tier.
        \item Speedups reported as bare means --- no sample- or seed-level
              dispersion.
        \item Rarely any statistical test against a matched autoregressive
              baseline.
      \end{itemize}
      \vspace{4pt}
      {\small This is a \emph{characterization} paper, not a new decoding
      algorithm.}
    \end{column}
    \begin{column}{0.54\textwidth}
      \begin{block}{Research questions}
        \small
        \begin{description}\itemsep2pt
          \item[RQ1] How much end-to-end speedup on consumer RTX hardware?
          \item[RQ2] How do $k$ and acceptance \emph{jointly} set net latency?
          \item[RQ3] How sensitive is the trade-off to deterministic
                vs.\ stochastic decoding?
          \item[RQ4] Do rankings and gains transfer RTX4090 $\to$
                RTX5090-laptop?
        \end{description}
      \end{block}
    \end{column}
  \end{columns}
\end{frame}
\note{%
\textbf{[2:00--2:40 | Gap and RQs]}\\[2pt]
Be explicit and unapologetic about the contribution type: ``we are not
proposing a faster drafter. We are asking whether the existing, unmodified
draft--verify loop delivers on the hardware most people actually own, and we
are answering it with the statistical care that deployment decisions
deserve.''\\[4pt]
Read RQ1 and RQ4 aloud; gesture at RQ2 and RQ3. RQ4 is the one the audience
will remember, so give it a beat: ``and the fourth question is the one that
surprised us.''\\[4pt]
\emph{If running long:} skip the three gaps, go straight to the RQ block ---
saves about 20 seconds.
}

% ================================================================
\section{Protocol}

\begin{frame}{Experimental protocol}
  \begin{columns}[T,onlytextwidth]
    \begin{column}{0.49\textwidth}
      \small
      \textbf{Models} (same-family pairings)
      \begin{itemize}\itemsep0pt
        \item Qwen2.5-3B-Instruct $\leftarrow$ 0.5B draft
        \item Qwen3-4B-Instruct $\leftarrow$ 0.6B draft
        \item FP16, matched tokenizers, KV cache
      \end{itemize}
      \vspace{3pt}
      \textbf{Hardware}
      \begin{itemize}\itemsep0pt
        \item RTX4090 desktop, 1008\,GB/s, 450\,W
        \item RTX5090 laptop, 896\,GB/s, $\leq$150\,W
      \end{itemize}
      \vspace{3pt}
      \textbf{Data} --- 1000 samples, frozen manifests:\\
      GSM8K 300 \, $\cdot$ \, MMLU 500 \, $\cdot$ \, CNN/DM 200
    \end{column}
    \begin{column}{0.49\textwidth}
      \begin{block}{Fixed-policy grid}
        \small
        $k \in \{4, 8, 16\}$ $\times$ \{deterministic, stochastic\}\\
        $\Rightarrow$ 6 configurations \textbf{per family per device},\\
        each against a \alert{regime-matched} AR baseline.
        \\[3pt]
        Deterministic: $T{=}0$, top-$p{=}1$.\quad
        Stochastic: $T{=}0.7$, top-$p{=}0.9$.
      \end{block}
      \begin{block}{Controls}
        \small
        Seed 42 primary; repeat seeds 123/999 at $k{=}4$.
        Frozen sample-ID manifests, identical prompts and length limits,
        warm-up before measurement, serial execution.
      \end{block}
    \end{column}
  \end{columns}
\end{frame}
\note{%
\textbf{[2:40--3:40 | Protocol]}\\[2pt]
The message here is \emph{pairing}. Every speculative run is compared to a
baseline that used the same prompts, the same sample IDs, the same length
limits and the same sampling regime --- so speedup can be computed per sample,
not just as a ratio of two independent averages. That is what makes the paired
$t$-tests later legitimate.\\[4pt]
Note the deliberate choice of same-family pairings: matching tokenizers means
no vocabulary-alignment overhead contaminating the measurement. Cross-family
pairing is a different study.\\[4pt]
On hardware, flag the counter-intuitive line: the laptop 5090 is the
\emph{newer} architecture --- Blackwell versus Ada --- but it has lower memory
bandwidth and roughly a third of the power budget. Plant that seed now; it pays
off at RQ4.\\[4pt]
Be honest about the seed budget: full grid at seed 42, repeats only at $k{=}4$.
Somebody will ask.
}

% ================================================================
\section{Results}

\begin{frame}{RQ1 --- On the RTX4090, it works}
  \centering
  \small
  \begin{tabular}{l cc cc cc}
    \toprule
    & \multicolumn{2}{c}{$k=4$} & \multicolumn{2}{c}{$k=8$}
    & \multicolumn{2}{c}{$k=16$}\\
    \cmidrule(lr){2-3}\cmidrule(lr){4-5}\cmidrule(lr){6-7}
    \textbf{Speedup $S$} & det & stoch & det & stoch & det & stoch\\
    \midrule
    Qwen2.5 (3B $\leftarrow$ 0.5B) & \best{3.0674} & 2.9394 & 2.4542 & 2.4063 & 1.8050 & 1.8185\\
    Qwen3 \; (4B $\leftarrow$ 0.6B) & \best{2.8548} & 2.4798 & 1.9648 & 1.6869 & 1.1091 & \bad{0.9719}\\
    \midrule
    Qwen3 acceptance $\alpha$ & 0.336 & 0.267 & 0.238 & 0.192 & 0.129 & 0.108\\
    Qwen3 block $\Beff$       & 1.33  & 1.06  & 1.86  & 1.50  & 1.96  & 1.64\\
    \bottomrule
  \end{tabular}
  \vspace{6pt}
  \begin{columns}[T]
    \begin{column}{0.48\textwidth}
      \small
      \begin{itemize}\itemsep1pt
        \item 11 of 12 configurations beat baseline.
        \item Both families peak at \best{deterministic $k{=}4$}.
      \end{itemize}
    \end{column}
    \begin{column}{0.48\textwidth}
      \small
      \begin{itemize}\itemsep1pt
        \item Qwen3 stochastic $k{=}16$ is the one loss (\bad{$0.97\times$}).
        \item Long proposals can erase the gain entirely.
      \end{itemize}
    \end{column}
  \end{columns}
\end{frame}
\note{%
\textbf{[3:40--4:55 | RQ1 headline]}\\[2pt]
This is the good-news slide --- deliver it confidently, then immediately
complicate it.\\[4pt]
Lead with the peak: ``on the RTX4090 we get just over three times on Qwen2.5
and just under three times on Qwen3, both at deterministic $k{=}4$. Eleven of
our twelve RTX4090 configurations beat their baseline.''\\[4pt]
Then the hook into RQ2: ``but notice which column the best number is in. It is
the \emph{smallest} $k$ we tested. Every step to a longer proposal costs us
speed. And in the bottom-right cell, Qwen3 stochastic at $k{=}16$, speculative
decoding is actually \emph{slower} than just running the target model
normally.''\\[4pt]
Also point at the two bottom rows now, because they set up the next slide:
acceptance $\alpha$ falls by roughly a factor of three from $k{=}4$ to $k{=}16$,
while the effective block size only rises from 1.3 to 2.0.\\[4pt]
\emph{Pointer discipline:} touch three cells only --- the two bold peaks and
the 0.9719. Do not read the table out.
}

% ================================================================
\begin{frame}{RQ2 --- Bigger blocks, slower generation}
  \begin{columns}[T,onlytextwidth]
    \begin{column}{0.55\textwidth}
      \panelA[\linewidth]
    \end{column}
    \begin{column}{0.43\textwidth}
      \small
      \textbf{Mean speedup, $k{=}4 \to k{=}16$}
      \begin{itemize}\itemsep0pt
        \item Qwen2.5: $2.9164 \to 1.8118$
        \item Qwen3: \; $2.6673 \to 1.0405$
      \end{itemize}
      \vspace{3pt}
      \textbf{Token timing degrades too}\\
      {\scriptsize (Qwen3, RTX4090, deterministic)}
      \begin{itemize}\itemsep0pt
        \item TTFT: \; $132.6 \to 349.6$\,ms
        \item TPOT: \; $41.2 \to 94.2$\,ms/token
      \end{itemize}
      \vspace{3pt}
      \begin{block}{}
        \small $\Beff$ \emph{rises} while latency \emph{worsens}: the
        overhead inflection arrives \alert{before $k{=}8$}.
      \end{block}
    \end{column}
  \end{columns}
\end{frame}
\note{%
\textbf{[4:55--5:55 | RQ2 the inflection]}\\[2pt]
This is the analytically interesting result, so slow down here.\\[4pt]
``Our hypothesis going in --- H2 --- was the textbook one: $k{=}8$ should beat
$k{=}4$, and $k{=}16$ might start to lose. That is wrong on this hardware. The
curve is monotonically decreasing from the smallest $k$ we tested. The
inflection has already happened before $k{=}8$.''\\[4pt]
Explain \emph{why} using the equation from earlier: the effective block size
does go up --- you really are committing more tokens per verify step --- but it
goes up sub-linearly while the drafting cost goes up linearly. Acceptance is
only 34\% at $k{=}4$ and collapses to 13\% at $k{=}16$, so most of those extra
$k$ draft steps are pure waste.\\[4pt]
The token-timing numbers on the right are the practitioner's punchline: at
$k{=}16$ you have nearly tripled time-to-first-token. For an interactive
application that is the worst possible place to spend latency --- the user sees
it directly.\\[4pt]
\emph{Implication to state:} on this class of hardware, tune $k$ downward, and
$k{=}4$ may not even be the floor.
}

% ================================================================
\begin{frame}{RQ3 --- Deterministic decoding is the safer default}
  \begin{columns}[T,onlytextwidth]
    \begin{column}{0.52\textwidth}
      \scriptsize
      \begin{tabular}{@{}lccc@{}}
        \toprule
        \textbf{Qwen3, RTX4090} & $k{=}4$ & $k{=}8$ & $k{=}16$\\
        \midrule
        Deterministic $S$ & 2.8548 & 1.9648 & 1.1091\\
        Stochastic $S$    & 2.4798 & 1.6869 & 0.9719\\
        \midrule
        Advantage         & \best{$+0.3750$} & $+0.2779$ & $+0.1372$\\
        \bottomrule
      \end{tabular}
      \vspace{8pt}
      \small
      \begin{itemize}\itemsep3pt
        \item Greedy drafting agrees with greedy verification more often
              $\Rightarrow$ higher $\alpha$ (0.336 vs.\ 0.267 at $k{=}4$).
        \item The advantage \emph{shrinks} with $k$ --- once acceptance is
              low, regime hardly matters.
      \end{itemize}
    \end{column}
    \begin{column}{0.46\textwidth}
      \begin{block}{Not universal}
        \small
        Qwen2.5 follows the same ordering at $k{=}4$ and $k{=}8$, but reverses
        slightly at $k{=}16$ ($1.8050$ det vs.\ $1.8185$ stoch).
      \end{block}
      \vspace{2pt}
      {\small\textbf{Takeaway.} If you need sampling, expect to pay
      $\sim$13\% of your speedup for it at low $k$.}
    \end{column}
  \end{columns}
\end{frame}
\note{%
\textbf{[5:55--6:30 | RQ3 regime]}\\[2pt]
Short slide --- this is the least surprising result, so move through it.\\[4pt]
The mechanism is the thing worth saying out loud: ``when both draft and target
decode greedily, the draft's proposal is compared against a deterministic
target continuation, so agreement is high. Turn sampling on for both and you
introduce two independent sources of randomness, acceptance drops from about
34\% to 27\%, and you lose roughly 13\% of your speedup.''\\[4pt]
Be careful to state the caveat rather than hiding it: at $k{=}16$ Qwen2.5
actually reverses, marginally. It is a small effect but we report it --- the
preference is a strong tendency, not a law.\\[4pt]
\emph{If running long:} this is the first slide to cut. State the one-line
takeaway --- ``deterministic wins at matched $k$, by a margin that shrinks as
$k$ grows'' --- and click through.
}

% ================================================================
\begin{frame}{RQ4 --- The same policy, a different GPU, and it all inverts}
  \begin{columns}[T,onlytextwidth]
    \begin{column}{0.53\textwidth}
      \panelB[\linewidth]
    \end{column}
    \begin{column}{0.45\textwidth}
      \small
      \textbf{Qwen3 on RTX5090-laptop}
      \scriptsize
      \begin{tabular}{@{}lcc@{}}
        \toprule
        Config & $S$ & $\alpha$\\
        \midrule
        $k{=}4$, det   & \bad{0.6743} & 0.335\\
        $k{=}4$, stoch & \bad{0.5924} & 0.285\\
        $k{=}8$, det   & \bad{0.4613} & 0.254\\
        $k{=}8$, stoch & \bad{0.3995} & 0.227\\
        $k{=}16$, det  & \bad{0.2588} & 0.180\\
        $k{=}16$, stoch& \bad{0.2222} & 0.154\\
        \bottomrule
      \end{tabular}
      \vspace{4pt}
      \small
      \begin{itemize}\itemsep1pt
        \item \alert{All six are slower than baseline.}
        \item Acceptance \emph{equals or beats} RTX4090 at every $k$ ---
              the loss is in execution cost, not the algorithm.
        \item Ordering survives: low $k$, deterministic.
      \end{itemize}
    \end{column}
  \end{columns}
\end{frame}
\note{%
\textbf{[6:30--7:50 | RQ4 portability --- the core result]}\\[2pt]
This is the slide the talk exists for. Pause before starting it.\\[4pt]
``We took the identical code, the identical models, the identical frozen
manifests, and moved to an RTX5090 laptop --- a \emph{newer} architecture. Every
single configuration is now slower than simply running the target model
autoregressively. The best case, deterministic $k{=}4$, gives us 0.67 --- we have
made inference one and a half times slower.''\\[4pt]
Then the diagnostic point, which is what makes this a finding rather than an
anecdote: ``acceptance is not the problem. At $k{=}4$ it is identical --- 0.336
against 0.335. At $k{=}8$ and $k{=}16$ the laptop actually accepts \emph{more}
--- 0.25 against 0.24, and 0.18 against 0.13. So the algorithm is doing at
least as well as on the desktop, and it is still slower everywhere. What
changed is the runtime economics: on this stack the draft model's per-step cost
relative to the target's verification cost is unfavourable, and speculative
decoding needs that ratio to be small.''\\[4pt]
Land the deployment rule: ``speedup is a property of the complete
model--hardware--runtime stack. It is not a property of the algorithm, and it is
certainly not a property of the GPU generation --- the newer card lost.''\\[4pt]
Finally, the constructive half: the \emph{ranking} did transfer. Low $k$ beats
high $k$, deterministic beats stochastic, on both devices. So policy
exploration on one machine is transferable; absolute claims are not.
}
\note{%
\textbf{[RQ4, continued]}\\[2pt]
\emph{Anticipate the question:} we do not claim to have isolated the cause. Our
candidate explanation is host-side orchestration and graph-capture behaviour ---
see the CUDA-graph note in the paper --- but we did not measure it, so we report
it as context, not as a cause. Say that before someone asks.\\[4pt]
\emph{Second likely question:} ``is the laptop card simply weaker?'' Partly ---
lower bandwidth and roughly a third of the power budget. But $S$ is a
\emph{ratio} against that same device's own autoregressive baseline, so a
uniformly slower card should not by itself push the ratio below one. The
backup hardware slide has the numbers.
}

% ================================================================
\begin{frame}{Is the effect real? Statistics and stability}
  \begin{columns}[T,onlytextwidth]
    \begin{column}{0.52\textwidth}
      \small
      \textbf{Paired one-sided $t$-tests}, $H_0{:}\ S \le 1$,\\
      paired by sample, Holm--Bonferroni across 6 configs
      (Qwen3, RTX4090, $n{=}1000$):
      \vspace{2pt}
      \scriptsize
      \begin{tabular}{@{}lcc@{}}
        \toprule
        Config & $d_S$ & $p_S$\\
        \midrule
        det $k{=}4,8,16$   & 1.18 / 1.10 / 0.29 & all $<10^{-18}$\\
        stoch $k{=}4,8$    & 1.24 / 1.13       & both $<10^{-179}$\\
        stoch $k{=}16$     & \bad{$-0.08$}     & \bad{$0.996$} \\
        \bottomrule
      \end{tabular}
      \vspace{6pt}
      \small
      \textbf{Dispersion is large.} Per-sample $S$ (det):
      $2.7561\pm0.5950$ at $k{=}4$ $\to$ $1.3549\pm0.6533$ at $k{=}16$.
      \\[2pt]
      {\scriptsize At $k{=}16$ the standard deviation is half the mean ---
      prompt-to-prompt variation dominates.}
    \end{column}
    \begin{column}{0.46\textwidth}
      \panelD[\linewidth]
      \vspace{-2pt}
      {\scriptsize Seed 999 vs.\ 123 at $k{=}4$: deterministic within
      $1.5\%$; stochastic $0.4\%$--$5.0\%$.}
    \end{column}
  \end{columns}
\end{frame}
\note{%
\textbf{[7:50--8:40 | Evidence quality]}\\[2pt]
Purpose of this slide: pre-empt the methodological objections, and show we are
not over-claiming.\\[4pt]
``Five of six RTX4090 configurations are significant after Holm--Bonferroni
correction, with large effect sizes. The sixth --- stochastic $k{=}16$ --- is
not significant, and its effect size is slightly negative, which is exactly
consistent with its sub-baseline mean. We report it as a null, not as noise.''\\[4pt]
The dispersion point matters more than the $p$-values: at $k{=}16$ the standard
deviation across prompts is roughly half the mean. A single reported average
speedup hides an enormous spread --- your slowest prompts may see no benefit at
all. This is one of the things the literature routinely omits.\\[4pt]
On the right panel: repeat seeds at $k{=}4$. Deterministic runs move by at most
1.5\% between seeds; stochastic by up to 5\%. Regime ordering is preserved
everywhere.\\[4pt]
\emph{Say the limitation before the reviewer does:} we only have repeat seeds at
$k{=}4$. We do not claim seed stability at $k{=}8$ or $k{=}16$.
}

% ================================================================
\begin{frame}{Speed is not free: quality on CNN/DailyMail}
  \centering
  \small
  \begin{tabular}{lccc}
    \toprule
    \textbf{Deterministic, $n{=}200$} & ROUGE-L & BLEU & BERTScore F1\\
    \midrule
    Autoregressive baseline & 19.58 & 3.20 & 0.793\\
    Speculative, $k{=}4$    & \bad{11.50} & \bad{1.39} & \bad{0.710}\\
    Speculative, $k{=}16$   & 14.92 & 2.04 & 0.741\\
    \bottomrule
  \end{tabular}
  \vspace{8pt}
  \begin{columns}[T]
    \begin{column}{0.47\textwidth}
      \small
      \begin{itemize}\itemsep2pt
        \item All three metrics agree: a real quality cost on the
              summarization task.
        \item It \alert{partially recovers} as $k$ grows --- moving
              \emph{against} the speed trend.
      \end{itemize}
    \end{column}
    \begin{column}{0.47\textwidth}
      \begin{block}{}
        \small The fastest configuration is also the weakest on
        summarization. There is a genuine speed--quality frontier here, not
        free acceleration.
      \end{block}
    \end{column}
  \end{columns}
\end{frame}
\note{%
\textbf{[8:40--9:25 | Quality trade-off]}\\[2pt]
Deliver this as an honest complication of our own headline number.\\[4pt]
``Earlier I said the peak was deterministic $k{=}4$. On summarization, that is
also our \emph{worst} quality point --- ROUGE-L drops from 19.6 to 11.5, and BLEU
and BERTScore move the same way. Three independent metrics agreeing means this
is not a metric artefact.''\\[4pt]
The direction is the interesting part: quality \emph{recovers} as $k$ increases,
which is the opposite direction to speed. So the two axes genuinely trade off
against each other across our grid, and the operating point is a choice, not a
default.\\[4pt]
\emph{Expected question --- be ready:} ``isn't speculative decoding supposed to
be distribution-preserving?'' Answer: exactly, and that is why this is worth
reporting. The guarantee is about the sampling distribution under the idealized
scheme; what we measure is an end-to-end implementation on a bounded generation
budget, where accepted-prefix commitment and length limits interact. We flag it
as an observed effect in our implementation and do not claim it refutes the
theory. Do not oversell this --- it is a limitation we are surfacing.\\[4pt]
Note that GSM8K and MMLU are accuracy tasks and are covered by our success
criteria; this slide is the summarization result specifically.
}

% ================================================================
\section{Wrap-up}

\begin{frame}{What we can and cannot claim}
  \begin{columns}[T,onlytextwidth]
    \begin{column}{0.48\textwidth}
      \textbf{Hypotheses}
      \small
      \begin{itemize}\itemsep2pt
        \item[\textbf{H1}] $S>1$ --- \best{supported} on RTX4090
              (11/12), \bad{rejected} on RTX5090-laptop.
        \item[\textbf{H2}] $k{=}8 > k{=}4$ --- \bad{rejected}; speedup
              falls monotonically.
        \item[\textbf{H3}] $\Beff$ up, latency worse --- \best{supported}.
        \item[\textbf{H4}] Deterministic faster --- \best{mostly supported}
              (one reversal).
      \end{itemize}
    \end{column}
    \begin{column}{0.48\textwidth}
      \textbf{Threats to validity}
      \small
      \begin{itemize}\itemsep2pt
        \item One PyTorch/Transformers stack; results do not extrapolate to
              other runtimes or GPUs.
        \item Same-family Qwen pairings only.
        \item No long-context, safety, or factual-consistency evaluation.
        \item Repeat seeds cover $k{=}4$ only.
        \item CUDA-graph capture differs across run families --- reported as
              context, \emph{not} as a measured cause.
      \end{itemize}
    \end{column}
  \end{columns}
\end{frame}
\note{%
\textbf{[9:25--9:55 | Scorecard and limits]}\\[2pt]
Move quickly --- this slide is mostly for the record and for the reviewers in
the room. Do not read every bullet.\\[4pt]
Highlight two things only. First, H2: ``we predicted $k{=}8$ would beat $k{=}4$
and it does not --- the overhead inflection is earlier than the folklore
suggests.'' Second, the limitations column exists because we want the claim
bounded: this is one runtime stack, one model family, one benchmark suite.\\[4pt]
The CUDA-graph bullet is the intellectually honest one --- we saw a difference in
capture behaviour between run families, and it is a plausible contributor to the
RTX5090 result, but we did not isolate it experimentally, so we refuse to call
it the cause.\\[4pt]
\emph{Timing checkpoint:} you should be at roughly 9:30 here. If you are past
10:30, skip straight to the takeaways.
}

% ================================================================
\begin{frame}[standout]
  \Large
  Policy ranking is portable. \\[6pt]
  Absolute speedup is not.
\end{frame}
\note{%
\textbf{[9:55--10:05 | The one line]}\\[2pt]
Pause. Let it sit for two seconds. This is the sentence you want people to
repeat to a colleague afterwards. Do not add anything to it --- click through
to the takeaways.
}

% ================================================================
\begin{frame}{Takeaways}
  \begin{enumerate}\itemsep6pt
    \item \textbf{Use a short proposal.} On both consumer GPUs and both model
          families, $k{=}4$ dominates $k{=}8$ and $k{=}16$ --- the overhead
          inflection arrives earlier than expected.
    \item \textbf{Prefer deterministic decoding} at matched $k$; sampling
          costs roughly 13\% of the speedup at $k{=}4$.
    \item \textbf{Benchmark on your target stack.} An RTX4090 gave up to
          $3.07\times$; the same code on an RTX5090 laptop was
          \emph{sub-baseline everywhere}, with unchanged acceptance.
    \item \textbf{Report dispersion and quality, not just a mean speedup.}
          Per-sample spread is large, and summarization quality moves against
          the speed trend.
  \end{enumerate}
  \vspace{6pt}
  \centering
  \small Frozen manifests, matched baselines, and raw per-sample artifacts
  support reproduction.
\end{frame}
\note{%
\textbf{[10:05--10:40 | Close]}\\[2pt]
Four sentences, one per point, then stop. Do not re-explain --- the audience has
seen all the evidence.\\[4pt]
Rule one: short proposals. Rule two: greedy if you can afford it. Rule three,
and this is the one to emphasize with your voice: measure on the machine you
will actually deploy on, because we have a concrete counter-example where a
newer GPU turned a three-times win into a one-and-a-half-times loss. Rule four:
report the spread and the quality, because the mean speedup on its own is not a
deployment decision.\\[4pt]
Close on future work in one breath: full grid across seeds, broader hardware and
model coverage, longer contexts and factuality-sensitive metrics.\\[4pt]
``Thank you --- I'm happy to take questions.''
}

% ================================================================
\appendix

\begin{frame}[standout]
  Backup slides
\end{frame}
\note{Backup material follows. Jump here only if asked.}

\begin{frame}{Backup --- full Qwen3 RTX4090 grid}
  \centering
  \small
  \begin{tabular}{lcccccc}
    \toprule
    \textbf{Config} & $S$ & $\alpha$ & $\Beff$ & $T_{\text{mean}}$ (s)
    & TTFT (ms) & TPOT (ms)\\
    \midrule
    $k{=}4$, det   & 2.8548 & 0.3362 & 1.33 & 4.8881 & 132.57 & 41.23\\
    $k{=}4$, stoch & 2.4798 & 0.2672 & 1.06 & 5.6196 & 134.54 & 45.51\\
    $k{=}8$, det   & 1.9648 & 0.2378 & 1.86 & 7.1024 & 204.99 & 58.76\\
    $k{=}8$, stoch & 1.6869 & 0.1916 & 1.50 & 8.2612 & 207.82 & 64.67\\
    $k{=}16$, det  & 1.1091 & 0.1292 & 1.96 & 12.5817 & 349.58 & 94.24\\
    $k{=}16$, stoch& 0.9719 & 0.1082 & 1.64 & 14.3390 & 354.82 & 106.03\\
    \bottomrule
  \end{tabular}
  \vspace{6pt}

  \small Target: Qwen3-4B-Instruct; draft: Qwen3-0.6B-Instruct; $n{=}1000$.
\end{frame}
\note{Use if asked for the full Qwen3 numbers, or for TTFT/TPOT detail.}

\begin{frame}{Backup --- full Qwen2.5 RTX4090 grid}
  \centering
  \small
  \begin{tabular}{lcccccc}
    \toprule
    \textbf{Config} & $S$ & $\alpha$ & $\Beff$ & $T_{\text{mean}}$ (s)
    & TTFT (ms) & TPOT (ms)\\
    \midrule
    $k{=}4$, det   & 3.0674 & 0.4212 & 1.66 & 3.7210 & 102.35 & 29.06\\
    $k{=}4$, stoch & 2.9394 & 0.4080 & 1.61 & 3.9631 & 105.26 & 30.89\\
    $k{=}8$, det   & 2.4542 & 0.3515 & 2.72 & 4.6508 & 162.76 & 36.73\\
    $k{=}8$, stoch & 2.4063 & 0.3444 & 2.67 & 4.8411 & 165.95 & 38.04\\
    $k{=}16$, det  & 1.8050 & 0.2539 & 3.76 & 6.3234 & 277.43 & 49.34\\
    $k{=}16$, stoch& 1.8185 & 0.2516 & 3.73 & 6.4059 & 278.85 & 49.81\\
    \bottomrule
  \end{tabular}
  \vspace{6pt}

  \small Target: Qwen2.5-3B-Instruct; draft: Qwen2.5-0.5B-Instruct; $n{=}1000$.\\
  Higher $\alpha$ and $\Beff$ than Qwen3 at every $k$.
\end{frame}
\note{%
Qwen2.5 accepts more than Qwen3 at every $k$ --- $\alpha=0.42$ vs.\ $0.34$ at
$k{=}4$, and effective block size 3.76 vs.\ 1.96 at $k{=}16$. That is why
Qwen2.5 never falls below baseline on the RTX4090 while Qwen3 does. Likely a
smaller capability gap between the 3B target and 0.5B draft than between the
4B target and 0.6B draft.
}

\begin{frame}{Backup --- acceptance portability}
  \centering
  \panelC[0.62\linewidth]
  \vspace{2pt}

  \small Qwen3 acceptance on the laptop is identical at $k{=}4$ (0.335 vs.\
  0.336) and \emph{higher} at $k{=}8$ and $k{=}16$ (0.254 vs.\ 0.238; 0.180
  vs.\ 0.129). The algorithm behaves at least as well; the
  \alert{runtime economics} differ.
\end{frame}
\note{%
The single most useful backup slide for the RQ4 discussion. If someone
challenges the RTX5090 result --- ``did you misconfigure it?'' --- this is the
answer: acceptance is a model-level quantity; it is unchanged at $k{=}4$ and
actually \emph{better} at higher $k$, so the speculative loop is doing at least
as well as it did on the desktop. The loss is in per-step execution cost, not
in the algorithm.
}

\begin{frame}{Backup --- success criteria and metrics}
  \begin{columns}[T,onlytextwidth]
    \begin{column}{0.48\textwidth}
      \small
      \textbf{Stated success criteria}
      \begin{itemize}\itemsep2pt
        \item Primary: mean $S \ge 1.3\times$ with Holm--Bonferroni-corrected
              $p_S < 0.05$.
        \item Primary: absolute quality drop $\le 1.0$ point on GSM8K and
              MMLU, deterministic regime.
        \item Secondary: best $(k, \text{regime})$ under quality constraints;
              dispersion; seed stability; cross-hardware ranking transfer.
      \end{itemize}
    \end{column}
    \begin{column}{0.48\textwidth}
      \small
      \textbf{Key metric definitions}
      \begin{itemize}\itemsep2pt
        \item $S$ = baseline latency / speculative latency
        \item $\alpha$ = accepted / proposed draft tokens
        \item $\Beff$ = accepted draft tokens / verify steps
        \item TTFT = first token time $-$ request start
        \item TPOT = (completion $-$ first token) / output tokens
        \item $d_S$ = Cohen's $d$, paired baseline vs.\ speculative latency
      \end{itemize}
    \end{column}
  \end{columns}
\end{frame}
\note{Reference slide for methodology questions.}

\begin{frame}{Backup --- hardware profiles}
  \centering
  \small
  \begin{tabular}{lll}
    \toprule
    \textbf{Component} & \textbf{RTX4090 desktop} & \textbf{RTX5090 laptop}\\
    \midrule
    GPU memory        & 24\,GB GDDR6X & 24\,GB GDDR7\\
    Memory bandwidth  & \best{1008\,GB/s} & 896\,GB/s\\
    CUDA cores        & \best{16384} & 10496\\
    Architecture      & Ada Lovelace & \best{Blackwell}\\
    Power profile     & \best{450\,W TGP class} & up to 150\,W TGP (OEM-dependent)\\
    Software stack    & \multicolumn{2}{c}{PyTorch + Transformers, pinned per run family}\\
    \bottomrule
  \end{tabular}
  \vspace{8pt}

  \small The laptop part is the newer architecture but has lower bandwidth and
  roughly \alert{one third of the power budget}.
\end{frame}
\note{%
The obvious question after RQ4 is ``is the 5090 laptop just a weaker card?''
Partly --- lower bandwidth, far lower sustained power. But note that this affects
the autoregressive baseline too, and $S$ is a \emph{ratio} against that same
device's baseline. So a uniformly slower card should not by itself flip the
ratio below 1. That is why we point at the draft/target cost ratio and the
runtime path rather than at raw card performance.
}

\end{document}
"
%    Dual screen :  set \def\NOTESMODE{} -> change option below
% ===============================================================
\documentclass[aspectratio=169, 11pt]{beamer}

% ---------------------------------------------------------------
% Shared preamble for the ICECCME 2026 presentation
% ---------------------------------------------------------------
\usetheme[progressbar=frametitle, numbering=fraction, block=fill]{metropolis}
\usepackage[T1]{fontenc}
\usepackage[utf8]{inputenc}
\usepackage{amsmath,amssymb}
\usepackage{booktabs}
\usepackage{graphicx}
\usepackage{xcolor}
\usepackage{appendixnumberbeamer}

% Colour-blind-safe palette, matching the paper figures.
\definecolor{cbBlue}{HTML}{0072B2}   % Qwen2.5 / RTX4090
\definecolor{cbOrange}{HTML}{D55E00} % Qwen3 / stochastic
\definecolor{cbGreen}{HTML}{009E73}  % RTX5090-laptop
\definecolor{cbGrey}{HTML}{555555}

\setbeamercolor{progress bar}{fg=cbBlue}
\setbeamercolor{alerted text}{fg=cbOrange}
\setbeamerfont{note page}{size=\small}

% Compact tables that survive a projector.
\setlength{\tabcolsep}{4.2pt}
\renewcommand{\arraystretch}{1.08}

% Shorthands used throughout the deck.
\newcommand{\up}{\textcolor{cbBlue}{$\uparrow$}}
\newcommand{\dn}{\textcolor{cbOrange}{$\downarrow$}}
\newcommand{\best}[1]{\textcolor{cbBlue}{\textbf{#1}}}
\newcommand{\bad}[1]{\textcolor{cbOrange}{\textbf{#1}}}
\newcommand{\Beff}{B_{\text{eff}}}

% Panel clipping for figures/specdec_configs_vs_baseline_metrics.pdf
% (a 2x2 matplotlib grid, 545.9 x 393.2 pt).
\newcommand{\panelA}[1][0.92\linewidth]{%
  \includegraphics[width=#1, trim=9 202 271 5, clip]%
  {../figures/specdec_configs_vs_baseline_metrics.pdf}}
\newcommand{\panelB}[1][0.92\linewidth]{%
  \includegraphics[width=#1, trim=271 202 22 5, clip]%
  {../figures/specdec_configs_vs_baseline_metrics.pdf}}
\newcommand{\panelC}[1][0.92\linewidth]{%
  \includegraphics[width=#1, trim=5 8 267 199, clip]%
  {../figures/specdec_configs_vs_baseline_metrics.pdf}}
\newcommand{\panelD}[1][0.92\linewidth]{%
  \includegraphics[width=#1, trim=271 8 5 199, clip]%
  {../figures/specdec_configs_vs_baseline_metrics.pdf}}


% --- notes switch -----------------------------------------------
\ifdefined\NOTESMODE
  \setbeameroption{show only notes}
  % Notes-only build: drop the slide thumbnail and shrink the type so that
  % even the longest note fits on a single page.
  \setbeamertemplate{note page}[plain]
  \setbeamerfont{note page}{size=\footnotesize}
\else
  \setbeameroption{hide notes}
  % For a presenter-display setup, comment the line above and use:
  % \setbeameroption{show notes on second screen=right}
\fi
% ----------------------------------------------------------------

\title{Characterizing Speculative Decoding Dynamics for Large
       Language Models on Consumer-Class GPUs}
\author{Samsu Sempena \and Zhuohan Cui}
\institute{The University of Sydney}
\date{ICECCME 2026 \\ 15--17 October 2026, Bali, Indonesia}

\begin{document}

% ================================================================
\begin{frame}[plain]
  \titlepage
\end{frame}
\note{%
\textbf{[0:00--0:20 | Title]}\\[2pt]
``Good morning. I'm Samsu Sempena from the University of Sydney, and this is
joint work with Zhuohan Cui. Our paper characterizes what speculative decoding
actually does on consumer-class GPUs --- the kind of hardware a practitioner
has under their desk, not in a data centre.''\\[4pt]
\emph{Cue:} say the one-line thesis up front so the audience has a hook ---
``the short version is that the \emph{policy ranking} travels between GPUs,
but the \emph{speedup} does not.''\\[4pt]
\emph{Housekeeping:} confirm the clicker works, and check whether the chair
wants questions during or after.
}

% ================================================================
\section{Motivation}

\begin{frame}{Decoding is sequential, and that is the bottleneck}
  \begin{columns}[T,onlytextwidth]
    \begin{column}{0.56\textwidth}
      \begin{itemize}
        \item Autoregressive decoding needs \alert{one full forward pass of the
              target model per token}.
        \item Single-stream, interactive latency is therefore bound by
              \emph{memory traffic}, not by arithmetic.
        \item Speculative decoding is the standard fix --- and it is reported
              to be a large, free win.
      \end{itemize}
      \vspace{4pt}
      \begin{block}{The practitioner's problem}
        Almost all published gains are measured on \textbf{data-centre}
        hardware. On a single RTX-class GPU, what actually happens?
      \end{block}
    \end{column}
    \begin{column}{0.42\textwidth}
      \scriptsize
      \begin{tabular}{@{}ll@{}}
        \toprule
        \multicolumn{2}{l}{\textbf{Open deployment questions}}\\
        \midrule
        Which draft size?      & 0.5B? 1.5B? \\
        Which proposal length? & $k=4$? $8$? $16$? \\
        Greedy or sampling?    & speed vs.\ quality \\
        Does it transfer?      & desktop $\to$ laptop \\
        \bottomrule
      \end{tabular}
    \end{column}
  \end{columns}
\end{frame}
\note{%
\textbf{[0:20--1:05 | Motivation]}\\[2pt]
Frame the bottleneck physically: every token costs a complete pass over the
target model's weights, so on a single stream you are moving gigabytes of
weights per token and the GPU's compute units are mostly idle. That is why
batch-1 interactive decoding is memory-bandwidth bound.\\[4pt]
Then the pivot: ``speculative decoding is the well-known answer, and the
literature reports two-to-five-times speedups. But look at where those numbers
come from --- A100s, H100s, B200s. The practitioner running a 3B or 4B model on
one RTX card has no reliable guidance.''\\[4pt]
Point at the four questions on the right --- these are exactly the knobs
someone has to set on day one, and say we will answer all four with
measurements.\\[4pt]
\emph{Do not} over-explain memory bandwidth; one sentence is enough.
}

% ================================================================
\begin{frame}{Speculative decoding in one slide}
  \begin{columns}[T,onlytextwidth]
    \begin{column}{0.50\textwidth}
      \textbf{Draft $\to$ verify $\to$ commit}
      \begin{enumerate}\itemsep2pt
        \item A small \textbf{draft} model proposes $k$ tokens
              autoregressively.
        \item The large \textbf{target} verifies all $k$ in
              \alert{one} parallel forward pass.
        \item Commit the accepted prefix; add one correction token on the
              first mismatch.
      \end{enumerate}
      \vspace{3pt}
      Modified rejection sampling keeps the output distribution
      \emph{identical} to the target's.
      \\[2pt]
      \tiny Leviathan et al.\ 2023; Chen et al.\ 2023
    \end{column}
    \begin{column}{0.48\textwidth}
      \begin{block}{Per-cycle latency}
        \[
          L \;=\; \frac{T_{\text{draft}} + T_{\text{verify}}}{\tau},
          \qquad
          T_{\text{draft}} = k \cdot t_{\text{draft step}}
        \]
        $\tau \in [1,\,k{+}1]$ accepted tokens per cycle.
      \end{block}
      \vspace{-2pt}
      \begin{itemize}\itemsep2pt
        \item Numerator grows \alert{linearly} in $k$.
        \item Denominator $\tau$ saturates as acceptance $\alpha$ falls.
        \item[$\Rightarrow$] There must be a $k$ beyond which it \emph{loses}.
      \end{itemize}
    \end{column}
  \end{columns}
\end{frame}
\note{%
\textbf{[1:05--2:00 | Method recap]}\\[2pt]
Keep this brisk --- most of the room knows the algorithm. Walk the three steps
once, then land hard on the equation, because it is the analytical spine of
every result that follows.\\[4pt]
Say it plainly: ``the cost in the numerator grows linearly with $k$, because
you run the draft model $k$ times. The payoff in the denominator is $\tau$, the
accepted tokens per cycle, and $\tau$ saturates --- once acceptance drops off,
extra proposals are wasted work. So the theory already tells you there is a
crossover. Our contribution is showing \emph{where} it sits on real consumer
hardware, and that it sits much lower than people assume.''\\[4pt]
Stress the correctness guarantee once --- ``this is lossless with respect to
the target distribution'' --- because it pre-empts a question, and because it
makes the quality result later more interesting.
}

% ================================================================
\begin{frame}{What is missing, and what we ask}
  \begin{columns}[T,onlytextwidth]
    \begin{column}{0.44\textwidth}
      \textbf{Three gaps in the literature}
      \begin{itemize}\itemsep3pt
        \item No cross-model \emph{and} cross-hardware validation at the
              consumer tier.
        \item Speedups reported as bare means --- no sample- or seed-level
              dispersion.
        \item Rarely any statistical test against a matched autoregressive
              baseline.
      \end{itemize}
      \vspace{4pt}
      {\small This is a \emph{characterization} paper, not a new decoding
      algorithm.}
    \end{column}
    \begin{column}{0.54\textwidth}
      \begin{block}{Research questions}
        \small
        \begin{description}\itemsep2pt
          \item[RQ1] How much end-to-end speedup on consumer RTX hardware?
          \item[RQ2] How do $k$ and acceptance \emph{jointly} set net latency?
          \item[RQ3] How sensitive is the trade-off to deterministic
                vs.\ stochastic decoding?
          \item[RQ4] Do rankings and gains transfer RTX4090 $\to$
                RTX5090-laptop?
        \end{description}
      \end{block}
    \end{column}
  \end{columns}
\end{frame}
\note{%
\textbf{[2:00--2:40 | Gap and RQs]}\\[2pt]
Be explicit and unapologetic about the contribution type: ``we are not
proposing a faster drafter. We are asking whether the existing, unmodified
draft--verify loop delivers on the hardware most people actually own, and we
are answering it with the statistical care that deployment decisions
deserve.''\\[4pt]
Read RQ1 and RQ4 aloud; gesture at RQ2 and RQ3. RQ4 is the one the audience
will remember, so give it a beat: ``and the fourth question is the one that
surprised us.''\\[4pt]
\emph{If running long:} skip the three gaps, go straight to the RQ block ---
saves about 20 seconds.
}

% ================================================================
\section{Protocol}

\begin{frame}{Experimental protocol}
  \begin{columns}[T,onlytextwidth]
    \begin{column}{0.49\textwidth}
      \small
      \textbf{Models} (same-family pairings)
      \begin{itemize}\itemsep0pt
        \item Qwen2.5-3B-Instruct $\leftarrow$ 0.5B draft
        \item Qwen3-4B-Instruct $\leftarrow$ 0.6B draft
        \item FP16, matched tokenizers, KV cache
      \end{itemize}
      \vspace{3pt}
      \textbf{Hardware}
      \begin{itemize}\itemsep0pt
        \item RTX4090 desktop, 1008\,GB/s, 450\,W
        \item RTX5090 laptop, 896\,GB/s, $\leq$150\,W
      \end{itemize}
      \vspace{3pt}
      \textbf{Data} --- 1000 samples, frozen manifests:\\
      GSM8K 300 \, $\cdot$ \, MMLU 500 \, $\cdot$ \, CNN/DM 200
    \end{column}
    \begin{column}{0.49\textwidth}
      \begin{block}{Fixed-policy grid}
        \small
        $k \in \{4, 8, 16\}$ $\times$ \{deterministic, stochastic\}\\
        $\Rightarrow$ 6 configurations \textbf{per family per device},\\
        each against a \alert{regime-matched} AR baseline.
        \\[3pt]
        Deterministic: $T{=}0$, top-$p{=}1$.\quad
        Stochastic: $T{=}0.7$, top-$p{=}0.9$.
      \end{block}
      \begin{block}{Controls}
        \small
        Seed 42 primary; repeat seeds 123/999 at $k{=}4$.
        Frozen sample-ID manifests, identical prompts and length limits,
        warm-up before measurement, serial execution.
      \end{block}
    \end{column}
  \end{columns}
\end{frame}
\note{%
\textbf{[2:40--3:40 | Protocol]}\\[2pt]
The message here is \emph{pairing}. Every speculative run is compared to a
baseline that used the same prompts, the same sample IDs, the same length
limits and the same sampling regime --- so speedup can be computed per sample,
not just as a ratio of two independent averages. That is what makes the paired
$t$-tests later legitimate.\\[4pt]
Note the deliberate choice of same-family pairings: matching tokenizers means
no vocabulary-alignment overhead contaminating the measurement. Cross-family
pairing is a different study.\\[4pt]
On hardware, flag the counter-intuitive line: the laptop 5090 is the
\emph{newer} architecture --- Blackwell versus Ada --- but it has lower memory
bandwidth and roughly a third of the power budget. Plant that seed now; it pays
off at RQ4.\\[4pt]
Be honest about the seed budget: full grid at seed 42, repeats only at $k{=}4$.
Somebody will ask.
}

% ================================================================
\section{Results}

\begin{frame}{RQ1 --- On the RTX4090, it works}
  \centering
  \small
  \begin{tabular}{l cc cc cc}
    \toprule
    & \multicolumn{2}{c}{$k=4$} & \multicolumn{2}{c}{$k=8$}
    & \multicolumn{2}{c}{$k=16$}\\
    \cmidrule(lr){2-3}\cmidrule(lr){4-5}\cmidrule(lr){6-7}
    \textbf{Speedup $S$} & det & stoch & det & stoch & det & stoch\\
    \midrule
    Qwen2.5 (3B $\leftarrow$ 0.5B) & \best{3.0674} & 2.9394 & 2.4542 & 2.4063 & 1.8050 & 1.8185\\
    Qwen3 \; (4B $\leftarrow$ 0.6B) & \best{2.8548} & 2.4798 & 1.9648 & 1.6869 & 1.1091 & \bad{0.9719}\\
    \midrule
    Qwen3 acceptance $\alpha$ & 0.336 & 0.267 & 0.238 & 0.192 & 0.129 & 0.108\\
    Qwen3 block $\Beff$       & 1.33  & 1.06  & 1.86  & 1.50  & 1.96  & 1.64\\
    \bottomrule
  \end{tabular}
  \vspace{6pt}
  \begin{columns}[T]
    \begin{column}{0.48\textwidth}
      \small
      \begin{itemize}\itemsep1pt
        \item 11 of 12 configurations beat baseline.
        \item Both families peak at \best{deterministic $k{=}4$}.
      \end{itemize}
    \end{column}
    \begin{column}{0.48\textwidth}
      \small
      \begin{itemize}\itemsep1pt
        \item Qwen3 stochastic $k{=}16$ is the one loss (\bad{$0.97\times$}).
        \item Long proposals can erase the gain entirely.
      \end{itemize}
    \end{column}
  \end{columns}
\end{frame}
\note{%
\textbf{[3:40--4:55 | RQ1 headline]}\\[2pt]
This is the good-news slide --- deliver it confidently, then immediately
complicate it.\\[4pt]
Lead with the peak: ``on the RTX4090 we get just over three times on Qwen2.5
and just under three times on Qwen3, both at deterministic $k{=}4$. Eleven of
our twelve RTX4090 configurations beat their baseline.''\\[4pt]
Then the hook into RQ2: ``but notice which column the best number is in. It is
the \emph{smallest} $k$ we tested. Every step to a longer proposal costs us
speed. And in the bottom-right cell, Qwen3 stochastic at $k{=}16$, speculative
decoding is actually \emph{slower} than just running the target model
normally.''\\[4pt]
Also point at the two bottom rows now, because they set up the next slide:
acceptance $\alpha$ falls by roughly a factor of three from $k{=}4$ to $k{=}16$,
while the effective block size only rises from 1.3 to 2.0.\\[4pt]
\emph{Pointer discipline:} touch three cells only --- the two bold peaks and
the 0.9719. Do not read the table out.
}

% ================================================================
\begin{frame}{RQ2 --- Bigger blocks, slower generation}
  \begin{columns}[T,onlytextwidth]
    \begin{column}{0.55\textwidth}
      \panelA[\linewidth]
    \end{column}
    \begin{column}{0.43\textwidth}
      \small
      \textbf{Mean speedup, $k{=}4 \to k{=}16$}
      \begin{itemize}\itemsep0pt
        \item Qwen2.5: $2.9164 \to 1.8118$
        \item Qwen3: \; $2.6673 \to 1.0405$
      \end{itemize}
      \vspace{3pt}
      \textbf{Token timing degrades too}\\
      {\scriptsize (Qwen3, RTX4090, deterministic)}
      \begin{itemize}\itemsep0pt
        \item TTFT: \; $132.6 \to 349.6$\,ms
        \item TPOT: \; $41.2 \to 94.2$\,ms/token
      \end{itemize}
      \vspace{3pt}
      \begin{block}{}
        \small $\Beff$ \emph{rises} while latency \emph{worsens}: the
        overhead inflection arrives \alert{before $k{=}8$}.
      \end{block}
    \end{column}
  \end{columns}
\end{frame}
\note{%
\textbf{[4:55--5:55 | RQ2 the inflection]}\\[2pt]
This is the analytically interesting result, so slow down here.\\[4pt]
``Our hypothesis going in --- H2 --- was the textbook one: $k{=}8$ should beat
$k{=}4$, and $k{=}16$ might start to lose. That is wrong on this hardware. The
curve is monotonically decreasing from the smallest $k$ we tested. The
inflection has already happened before $k{=}8$.''\\[4pt]
Explain \emph{why} using the equation from earlier: the effective block size
does go up --- you really are committing more tokens per verify step --- but it
goes up sub-linearly while the drafting cost goes up linearly. Acceptance is
only 34\% at $k{=}4$ and collapses to 13\% at $k{=}16$, so most of those extra
$k$ draft steps are pure waste.\\[4pt]
The token-timing numbers on the right are the practitioner's punchline: at
$k{=}16$ you have nearly tripled time-to-first-token. For an interactive
application that is the worst possible place to spend latency --- the user sees
it directly.\\[4pt]
\emph{Implication to state:} on this class of hardware, tune $k$ downward, and
$k{=}4$ may not even be the floor.
}

% ================================================================
\begin{frame}{RQ3 --- Deterministic decoding is the safer default}
  \begin{columns}[T,onlytextwidth]
    \begin{column}{0.52\textwidth}
      \scriptsize
      \begin{tabular}{@{}lccc@{}}
        \toprule
        \textbf{Qwen3, RTX4090} & $k{=}4$ & $k{=}8$ & $k{=}16$\\
        \midrule
        Deterministic $S$ & 2.8548 & 1.9648 & 1.1091\\
        Stochastic $S$    & 2.4798 & 1.6869 & 0.9719\\
        \midrule
        Advantage         & \best{$+0.3750$} & $+0.2779$ & $+0.1372$\\
        \bottomrule
      \end{tabular}
      \vspace{8pt}
      \small
      \begin{itemize}\itemsep3pt
        \item Greedy drafting agrees with greedy verification more often
              $\Rightarrow$ higher $\alpha$ (0.336 vs.\ 0.267 at $k{=}4$).
        \item The advantage \emph{shrinks} with $k$ --- once acceptance is
              low, regime hardly matters.
      \end{itemize}
    \end{column}
    \begin{column}{0.46\textwidth}
      \begin{block}{Not universal}
        \small
        Qwen2.5 follows the same ordering at $k{=}4$ and $k{=}8$, but reverses
        slightly at $k{=}16$ ($1.8050$ det vs.\ $1.8185$ stoch).
      \end{block}
      \vspace{2pt}
      {\small\textbf{Takeaway.} If you need sampling, expect to pay
      $\sim$13\% of your speedup for it at low $k$.}
    \end{column}
  \end{columns}
\end{frame}
\note{%
\textbf{[5:55--6:30 | RQ3 regime]}\\[2pt]
Short slide --- this is the least surprising result, so move through it.\\[4pt]
The mechanism is the thing worth saying out loud: ``when both draft and target
decode greedily, the draft's proposal is compared against a deterministic
target continuation, so agreement is high. Turn sampling on for both and you
introduce two independent sources of randomness, acceptance drops from about
34\% to 27\%, and you lose roughly 13\% of your speedup.''\\[4pt]
Be careful to state the caveat rather than hiding it: at $k{=}16$ Qwen2.5
actually reverses, marginally. It is a small effect but we report it --- the
preference is a strong tendency, not a law.\\[4pt]
\emph{If running long:} this is the first slide to cut. State the one-line
takeaway --- ``deterministic wins at matched $k$, by a margin that shrinks as
$k$ grows'' --- and click through.
}

% ================================================================
\begin{frame}{RQ4 --- The same policy, a different GPU, and it all inverts}
  \begin{columns}[T,onlytextwidth]
    \begin{column}{0.53\textwidth}
      \panelB[\linewidth]
    \end{column}
    \begin{column}{0.45\textwidth}
      \small
      \textbf{Qwen3 on RTX5090-laptop}
      \scriptsize
      \begin{tabular}{@{}lcc@{}}
        \toprule
        Config & $S$ & $\alpha$\\
        \midrule
        $k{=}4$, det   & \bad{0.6743} & 0.335\\
        $k{=}4$, stoch & \bad{0.5924} & 0.285\\
        $k{=}8$, det   & \bad{0.4613} & 0.254\\
        $k{=}8$, stoch & \bad{0.3995} & 0.227\\
        $k{=}16$, det  & \bad{0.2588} & 0.180\\
        $k{=}16$, stoch& \bad{0.2222} & 0.154\\
        \bottomrule
      \end{tabular}
      \vspace{4pt}
      \small
      \begin{itemize}\itemsep1pt
        \item \alert{All six are slower than baseline.}
        \item Acceptance \emph{equals or beats} RTX4090 at every $k$ ---
              the loss is in execution cost, not the algorithm.
        \item Ordering survives: low $k$, deterministic.
      \end{itemize}
    \end{column}
  \end{columns}
\end{frame}
\note{%
\textbf{[6:30--7:50 | RQ4 portability --- the core result]}\\[2pt]
This is the slide the talk exists for. Pause before starting it.\\[4pt]
``We took the identical code, the identical models, the identical frozen
manifests, and moved to an RTX5090 laptop --- a \emph{newer} architecture. Every
single configuration is now slower than simply running the target model
autoregressively. The best case, deterministic $k{=}4$, gives us 0.67 --- we have
made inference one and a half times slower.''\\[4pt]
Then the diagnostic point, which is what makes this a finding rather than an
anecdote: ``acceptance is not the problem. At $k{=}4$ it is identical --- 0.336
against 0.335. At $k{=}8$ and $k{=}16$ the laptop actually accepts \emph{more}
--- 0.25 against 0.24, and 0.18 against 0.13. So the algorithm is doing at
least as well as on the desktop, and it is still slower everywhere. What
changed is the runtime economics: on this stack the draft model's per-step cost
relative to the target's verification cost is unfavourable, and speculative
decoding needs that ratio to be small.''\\[4pt]
Land the deployment rule: ``speedup is a property of the complete
model--hardware--runtime stack. It is not a property of the algorithm, and it is
certainly not a property of the GPU generation --- the newer card lost.''\\[4pt]
Finally, the constructive half: the \emph{ranking} did transfer. Low $k$ beats
high $k$, deterministic beats stochastic, on both devices. So policy
exploration on one machine is transferable; absolute claims are not.
}
\note{%
\textbf{[RQ4, continued]}\\[2pt]
\emph{Anticipate the question:} we do not claim to have isolated the cause. Our
candidate explanation is host-side orchestration and graph-capture behaviour ---
see the CUDA-graph note in the paper --- but we did not measure it, so we report
it as context, not as a cause. Say that before someone asks.\\[4pt]
\emph{Second likely question:} ``is the laptop card simply weaker?'' Partly ---
lower bandwidth and roughly a third of the power budget. But $S$ is a
\emph{ratio} against that same device's own autoregressive baseline, so a
uniformly slower card should not by itself push the ratio below one. The
backup hardware slide has the numbers.
}

% ================================================================
\begin{frame}{Is the effect real? Statistics and stability}
  \begin{columns}[T,onlytextwidth]
    \begin{column}{0.52\textwidth}
      \small
      \textbf{Paired one-sided $t$-tests}, $H_0{:}\ S \le 1$,\\
      paired by sample, Holm--Bonferroni across 6 configs
      (Qwen3, RTX4090, $n{=}1000$):
      \vspace{2pt}
      \scriptsize
      \begin{tabular}{@{}lcc@{}}
        \toprule
        Config & $d_S$ & $p_S$\\
        \midrule
        det $k{=}4,8,16$   & 1.18 / 1.10 / 0.29 & all $<10^{-18}$\\
        stoch $k{=}4,8$    & 1.24 / 1.13       & both $<10^{-179}$\\
        stoch $k{=}16$     & \bad{$-0.08$}     & \bad{$0.996$} \\
        \bottomrule
      \end{tabular}
      \vspace{6pt}
      \small
      \textbf{Dispersion is large.} Per-sample $S$ (det):
      $2.7561\pm0.5950$ at $k{=}4$ $\to$ $1.3549\pm0.6533$ at $k{=}16$.
      \\[2pt]
      {\scriptsize At $k{=}16$ the standard deviation is half the mean ---
      prompt-to-prompt variation dominates.}
    \end{column}
    \begin{column}{0.46\textwidth}
      \panelD[\linewidth]
      \vspace{-2pt}
      {\scriptsize Seed 999 vs.\ 123 at $k{=}4$: deterministic within
      $1.5\%$; stochastic $0.4\%$--$5.0\%$.}
    \end{column}
  \end{columns}
\end{frame}
\note{%
\textbf{[7:50--8:40 | Evidence quality]}\\[2pt]
Purpose of this slide: pre-empt the methodological objections, and show we are
not over-claiming.\\[4pt]
``Five of six RTX4090 configurations are significant after Holm--Bonferroni
correction, with large effect sizes. The sixth --- stochastic $k{=}16$ --- is
not significant, and its effect size is slightly negative, which is exactly
consistent with its sub-baseline mean. We report it as a null, not as noise.''\\[4pt]
The dispersion point matters more than the $p$-values: at $k{=}16$ the standard
deviation across prompts is roughly half the mean. A single reported average
speedup hides an enormous spread --- your slowest prompts may see no benefit at
all. This is one of the things the literature routinely omits.\\[4pt]
On the right panel: repeat seeds at $k{=}4$. Deterministic runs move by at most
1.5\% between seeds; stochastic by up to 5\%. Regime ordering is preserved
everywhere.\\[4pt]
\emph{Say the limitation before the reviewer does:} we only have repeat seeds at
$k{=}4$. We do not claim seed stability at $k{=}8$ or $k{=}16$.
}

% ================================================================
\begin{frame}{Speed is not free: quality on CNN/DailyMail}
  \centering
  \small
  \begin{tabular}{lccc}
    \toprule
    \textbf{Deterministic, $n{=}200$} & ROUGE-L & BLEU & BERTScore F1\\
    \midrule
    Autoregressive baseline & 19.58 & 3.20 & 0.793\\
    Speculative, $k{=}4$    & \bad{11.50} & \bad{1.39} & \bad{0.710}\\
    Speculative, $k{=}16$   & 14.92 & 2.04 & 0.741\\
    \bottomrule
  \end{tabular}
  \vspace{8pt}
  \begin{columns}[T]
    \begin{column}{0.47\textwidth}
      \small
      \begin{itemize}\itemsep2pt
        \item All three metrics agree: a real quality cost on the
              summarization task.
        \item It \alert{partially recovers} as $k$ grows --- moving
              \emph{against} the speed trend.
      \end{itemize}
    \end{column}
    \begin{column}{0.47\textwidth}
      \begin{block}{}
        \small The fastest configuration is also the weakest on
        summarization. There is a genuine speed--quality frontier here, not
        free acceleration.
      \end{block}
    \end{column}
  \end{columns}
\end{frame}
\note{%
\textbf{[8:40--9:25 | Quality trade-off]}\\[2pt]
Deliver this as an honest complication of our own headline number.\\[4pt]
``Earlier I said the peak was deterministic $k{=}4$. On summarization, that is
also our \emph{worst} quality point --- ROUGE-L drops from 19.6 to 11.5, and BLEU
and BERTScore move the same way. Three independent metrics agreeing means this
is not a metric artefact.''\\[4pt]
The direction is the interesting part: quality \emph{recovers} as $k$ increases,
which is the opposite direction to speed. So the two axes genuinely trade off
against each other across our grid, and the operating point is a choice, not a
default.\\[4pt]
\emph{Expected question --- be ready:} ``isn't speculative decoding supposed to
be distribution-preserving?'' Answer: exactly, and that is why this is worth
reporting. The guarantee is about the sampling distribution under the idealized
scheme; what we measure is an end-to-end implementation on a bounded generation
budget, where accepted-prefix commitment and length limits interact. We flag it
as an observed effect in our implementation and do not claim it refutes the
theory. Do not oversell this --- it is a limitation we are surfacing.\\[4pt]
Note that GSM8K and MMLU are accuracy tasks and are covered by our success
criteria; this slide is the summarization result specifically.
}

% ================================================================
\section{Wrap-up}

\begin{frame}{What we can and cannot claim}
  \begin{columns}[T,onlytextwidth]
    \begin{column}{0.48\textwidth}
      \textbf{Hypotheses}
      \small
      \begin{itemize}\itemsep2pt
        \item[\textbf{H1}] $S>1$ --- \best{supported} on RTX4090
              (11/12), \bad{rejected} on RTX5090-laptop.
        \item[\textbf{H2}] $k{=}8 > k{=}4$ --- \bad{rejected}; speedup
              falls monotonically.
        \item[\textbf{H3}] $\Beff$ up, latency worse --- \best{supported}.
        \item[\textbf{H4}] Deterministic faster --- \best{mostly supported}
              (one reversal).
      \end{itemize}
    \end{column}
    \begin{column}{0.48\textwidth}
      \textbf{Threats to validity}
      \small
      \begin{itemize}\itemsep2pt
        \item One PyTorch/Transformers stack; results do not extrapolate to
              other runtimes or GPUs.
        \item Same-family Qwen pairings only.
        \item No long-context, safety, or factual-consistency evaluation.
        \item Repeat seeds cover $k{=}4$ only.
        \item CUDA-graph capture differs across run families --- reported as
              context, \emph{not} as a measured cause.
      \end{itemize}
    \end{column}
  \end{columns}
\end{frame}
\note{%
\textbf{[9:25--9:55 | Scorecard and limits]}\\[2pt]
Move quickly --- this slide is mostly for the record and for the reviewers in
the room. Do not read every bullet.\\[4pt]
Highlight two things only. First, H2: ``we predicted $k{=}8$ would beat $k{=}4$
and it does not --- the overhead inflection is earlier than the folklore
suggests.'' Second, the limitations column exists because we want the claim
bounded: this is one runtime stack, one model family, one benchmark suite.\\[4pt]
The CUDA-graph bullet is the intellectually honest one --- we saw a difference in
capture behaviour between run families, and it is a plausible contributor to the
RTX5090 result, but we did not isolate it experimentally, so we refuse to call
it the cause.\\[4pt]
\emph{Timing checkpoint:} you should be at roughly 9:30 here. If you are past
10:30, skip straight to the takeaways.
}

% ================================================================
\begin{frame}[standout]
  \Large
  Policy ranking is portable. \\[6pt]
  Absolute speedup is not.
\end{frame}
\note{%
\textbf{[9:55--10:05 | The one line]}\\[2pt]
Pause. Let it sit for two seconds. This is the sentence you want people to
repeat to a colleague afterwards. Do not add anything to it --- click through
to the takeaways.
}

% ================================================================
\begin{frame}{Takeaways}
  \begin{enumerate}\itemsep6pt
    \item \textbf{Use a short proposal.} On both consumer GPUs and both model
          families, $k{=}4$ dominates $k{=}8$ and $k{=}16$ --- the overhead
          inflection arrives earlier than expected.
    \item \textbf{Prefer deterministic decoding} at matched $k$; sampling
          costs roughly 13\% of the speedup at $k{=}4$.
    \item \textbf{Benchmark on your target stack.} An RTX4090 gave up to
          $3.07\times$; the same code on an RTX5090 laptop was
          \emph{sub-baseline everywhere}, with unchanged acceptance.
    \item \textbf{Report dispersion and quality, not just a mean speedup.}
          Per-sample spread is large, and summarization quality moves against
          the speed trend.
  \end{enumerate}
  \vspace{6pt}
  \centering
  \small Frozen manifests, matched baselines, and raw per-sample artifacts
  support reproduction.
\end{frame}
\note{%
\textbf{[10:05--10:40 | Close]}\\[2pt]
Four sentences, one per point, then stop. Do not re-explain --- the audience has
seen all the evidence.\\[4pt]
Rule one: short proposals. Rule two: greedy if you can afford it. Rule three,
and this is the one to emphasize with your voice: measure on the machine you
will actually deploy on, because we have a concrete counter-example where a
newer GPU turned a three-times win into a one-and-a-half-times loss. Rule four:
report the spread and the quality, because the mean speedup on its own is not a
deployment decision.\\[4pt]
Close on future work in one breath: full grid across seeds, broader hardware and
model coverage, longer contexts and factuality-sensitive metrics.\\[4pt]
``Thank you --- I'm happy to take questions.''
}

% ================================================================
\appendix

\begin{frame}[standout]
  Backup slides
\end{frame}
\note{Backup material follows. Jump here only if asked.}

\begin{frame}{Backup --- full Qwen3 RTX4090 grid}
  \centering
  \small
  \begin{tabular}{lcccccc}
    \toprule
    \textbf{Config} & $S$ & $\alpha$ & $\Beff$ & $T_{\text{mean}}$ (s)
    & TTFT (ms) & TPOT (ms)\\
    \midrule
    $k{=}4$, det   & 2.8548 & 0.3362 & 1.33 & 4.8881 & 132.57 & 41.23\\
    $k{=}4$, stoch & 2.4798 & 0.2672 & 1.06 & 5.6196 & 134.54 & 45.51\\
    $k{=}8$, det   & 1.9648 & 0.2378 & 1.86 & 7.1024 & 204.99 & 58.76\\
    $k{=}8$, stoch & 1.6869 & 0.1916 & 1.50 & 8.2612 & 207.82 & 64.67\\
    $k{=}16$, det  & 1.1091 & 0.1292 & 1.96 & 12.5817 & 349.58 & 94.24\\
    $k{=}16$, stoch& 0.9719 & 0.1082 & 1.64 & 14.3390 & 354.82 & 106.03\\
    \bottomrule
  \end{tabular}
  \vspace{6pt}

  \small Target: Qwen3-4B-Instruct; draft: Qwen3-0.6B-Instruct; $n{=}1000$.
\end{frame}
\note{Use if asked for the full Qwen3 numbers, or for TTFT/TPOT detail.}

\begin{frame}{Backup --- full Qwen2.5 RTX4090 grid}
  \centering
  \small
  \begin{tabular}{lcccccc}
    \toprule
    \textbf{Config} & $S$ & $\alpha$ & $\Beff$ & $T_{\text{mean}}$ (s)
    & TTFT (ms) & TPOT (ms)\\
    \midrule
    $k{=}4$, det   & 3.0674 & 0.4212 & 1.66 & 3.7210 & 102.35 & 29.06\\
    $k{=}4$, stoch & 2.9394 & 0.4080 & 1.61 & 3.9631 & 105.26 & 30.89\\
    $k{=}8$, det   & 2.4542 & 0.3515 & 2.72 & 4.6508 & 162.76 & 36.73\\
    $k{=}8$, stoch & 2.4063 & 0.3444 & 2.67 & 4.8411 & 165.95 & 38.04\\
    $k{=}16$, det  & 1.8050 & 0.2539 & 3.76 & 6.3234 & 277.43 & 49.34\\
    $k{=}16$, stoch& 1.8185 & 0.2516 & 3.73 & 6.4059 & 278.85 & 49.81\\
    \bottomrule
  \end{tabular}
  \vspace{6pt}

  \small Target: Qwen2.5-3B-Instruct; draft: Qwen2.5-0.5B-Instruct; $n{=}1000$.\\
  Higher $\alpha$ and $\Beff$ than Qwen3 at every $k$.
\end{frame}
\note{%
Qwen2.5 accepts more than Qwen3 at every $k$ --- $\alpha=0.42$ vs.\ $0.34$ at
$k{=}4$, and effective block size 3.76 vs.\ 1.96 at $k{=}16$. That is why
Qwen2.5 never falls below baseline on the RTX4090 while Qwen3 does. Likely a
smaller capability gap between the 3B target and 0.5B draft than between the
4B target and 0.6B draft.
}

\begin{frame}{Backup --- acceptance portability}
  \centering
  \panelC[0.62\linewidth]
  \vspace{2pt}

  \small Qwen3 acceptance on the laptop is identical at $k{=}4$ (0.335 vs.\
  0.336) and \emph{higher} at $k{=}8$ and $k{=}16$ (0.254 vs.\ 0.238; 0.180
  vs.\ 0.129). The algorithm behaves at least as well; the
  \alert{runtime economics} differ.
\end{frame}
\note{%
The single most useful backup slide for the RQ4 discussion. If someone
challenges the RTX5090 result --- ``did you misconfigure it?'' --- this is the
answer: acceptance is a model-level quantity; it is unchanged at $k{=}4$ and
actually \emph{better} at higher $k$, so the speculative loop is doing at least
as well as it did on the desktop. The loss is in per-step execution cost, not
in the algorithm.
}

\begin{frame}{Backup --- success criteria and metrics}
  \begin{columns}[T,onlytextwidth]
    \begin{column}{0.48\textwidth}
      \small
      \textbf{Stated success criteria}
      \begin{itemize}\itemsep2pt
        \item Primary: mean $S \ge 1.3\times$ with Holm--Bonferroni-corrected
              $p_S < 0.05$.
        \item Primary: absolute quality drop $\le 1.0$ point on GSM8K and
              MMLU, deterministic regime.
        \item Secondary: best $(k, \text{regime})$ under quality constraints;
              dispersion; seed stability; cross-hardware ranking transfer.
      \end{itemize}
    \end{column}
    \begin{column}{0.48\textwidth}
      \small
      \textbf{Key metric definitions}
      \begin{itemize}\itemsep2pt
        \item $S$ = baseline latency / speculative latency
        \item $\alpha$ = accepted / proposed draft tokens
        \item $\Beff$ = accepted draft tokens / verify steps
        \item TTFT = first token time $-$ request start
        \item TPOT = (completion $-$ first token) / output tokens
        \item $d_S$ = Cohen's $d$, paired baseline vs.\ speculative latency
      \end{itemize}
    \end{column}
  \end{columns}
\end{frame}
\note{Reference slide for methodology questions.}

\begin{frame}{Backup --- hardware profiles}
  \centering
  \small
  \begin{tabular}{lll}
    \toprule
    \textbf{Component} & \textbf{RTX4090 desktop} & \textbf{RTX5090 laptop}\\
    \midrule
    GPU memory        & 24\,GB GDDR6X & 24\,GB GDDR7\\
    Memory bandwidth  & \best{1008\,GB/s} & 896\,GB/s\\
    CUDA cores        & \best{16384} & 10496\\
    Architecture      & Ada Lovelace & \best{Blackwell}\\
    Power profile     & \best{450\,W TGP class} & up to 150\,W TGP (OEM-dependent)\\
    Software stack    & \multicolumn{2}{c}{PyTorch + Transformers, pinned per run family}\\
    \bottomrule
  \end{tabular}
  \vspace{8pt}

  \small The laptop part is the newer architecture but has lower bandwidth and
  roughly \alert{one third of the power budget}.
\end{frame}
\note{%
The obvious question after RQ4 is ``is the 5090 laptop just a weaker card?''
Partly --- lower bandwidth, far lower sustained power. But note that this affects
the autoregressive baseline too, and $S$ is a \emph{ratio} against that same
device's baseline. So a uniformly slower card should not by itself flip the
ratio below 1. That is why we point at the draft/target cost ratio and the
runtime path rather than at raw card performance.
}

\end{document}
"
%    Dual screen :  set \def\NOTESMODE{} -> change option below
% ===============================================================
\documentclass[aspectratio=169, 11pt]{beamer}

% ---------------------------------------------------------------
% Shared preamble for the ICECCME 2026 presentation
% ---------------------------------------------------------------
\usetheme[progressbar=frametitle, numbering=fraction, block=fill]{metropolis}
\usepackage[T1]{fontenc}
\usepackage[utf8]{inputenc}
\usepackage{amsmath,amssymb}
\usepackage{booktabs}
\usepackage{graphicx}
\usepackage{xcolor}
\usepackage{appendixnumberbeamer}

% Colour-blind-safe palette, matching the paper figures.
\definecolor{cbBlue}{HTML}{0072B2}   % Qwen2.5 / RTX4090
\definecolor{cbOrange}{HTML}{D55E00} % Qwen3 / stochastic
\definecolor{cbGreen}{HTML}{009E73}  % RTX5090-laptop
\definecolor{cbGrey}{HTML}{555555}

\setbeamercolor{progress bar}{fg=cbBlue}
\setbeamercolor{alerted text}{fg=cbOrange}
\setbeamerfont{note page}{size=\small}

% Compact tables that survive a projector.
\setlength{\tabcolsep}{4.2pt}
\renewcommand{\arraystretch}{1.08}

% Shorthands used throughout the deck.
\newcommand{\up}{\textcolor{cbBlue}{$\uparrow$}}
\newcommand{\dn}{\textcolor{cbOrange}{$\downarrow$}}
\newcommand{\best}[1]{\textcolor{cbBlue}{\textbf{#1}}}
\newcommand{\bad}[1]{\textcolor{cbOrange}{\textbf{#1}}}
\newcommand{\Beff}{B_{\text{eff}}}

% Panel clipping for figures/specdec_configs_vs_baseline_metrics.pdf
% (a 2x2 matplotlib grid, 545.9 x 393.2 pt).
\newcommand{\panelA}[1][0.92\linewidth]{%
  \includegraphics[width=#1, trim=9 202 271 5, clip]%
  {../figures/specdec_configs_vs_baseline_metrics.pdf}}
\newcommand{\panelB}[1][0.92\linewidth]{%
  \includegraphics[width=#1, trim=271 202 22 5, clip]%
  {../figures/specdec_configs_vs_baseline_metrics.pdf}}
\newcommand{\panelC}[1][0.92\linewidth]{%
  \includegraphics[width=#1, trim=5 8 267 199, clip]%
  {../figures/specdec_configs_vs_baseline_metrics.pdf}}
\newcommand{\panelD}[1][0.92\linewidth]{%
  \includegraphics[width=#1, trim=271 8 5 199, clip]%
  {../figures/specdec_configs_vs_baseline_metrics.pdf}}


% --- notes switch -----------------------------------------------
\ifdefined\NOTESMODE
  \setbeameroption{show only notes}
  % Notes-only build: drop the slide thumbnail and shrink the type so that
  % even the longest note fits on a single page.
  \setbeamertemplate{note page}[plain]
  \setbeamerfont{note page}{size=\footnotesize}
\else
  \setbeameroption{hide notes}
  % For a presenter-display setup, comment the line above and use:
  % \setbeameroption{show notes on second screen=right}
\fi
% ----------------------------------------------------------------

\title{Characterizing Speculative Decoding Dynamics for Large
       Language Models on Consumer-Class GPUs}
\author{Samsu Sempena \and Zhuohan Cui}
\institute{The University of Sydney}
\date{ICECCME 2026 \\ 15--17 October 2026, Bali, Indonesia}

\begin{document}

% ================================================================
\begin{frame}[plain]
  \titlepage
\end{frame}
\note{%
\textbf{[0:00--0:20 | Title]}\\[2pt]
``Good morning. I'm Samsu Sempena from the University of Sydney, and this is
joint work with Zhuohan Cui. Our paper characterizes what speculative decoding
actually does on consumer-class GPUs --- the kind of hardware a practitioner
has under their desk, not in a data centre.''\\[4pt]
\emph{Cue:} say the one-line thesis up front so the audience has a hook ---
``the short version is that the \emph{policy ranking} travels between GPUs,
but the \emph{speedup} does not.''\\[4pt]
\emph{Housekeeping:} confirm the clicker works, and check whether the chair
wants questions during or after.
}

% ================================================================
\section{Motivation}

\begin{frame}{Decoding is sequential, and that is the bottleneck}
  \begin{columns}[T,onlytextwidth]
    \begin{column}{0.56\textwidth}
      \begin{itemize}
        \item Autoregressive decoding needs \alert{one full forward pass of the
              target model per token}.
        \item Single-stream, interactive latency is therefore bound by
              \emph{memory traffic}, not by arithmetic.
        \item Speculative decoding is the standard fix --- and it is reported
              to be a large, free win.
      \end{itemize}
      \vspace{4pt}
      \begin{block}{The practitioner's problem}
        Almost all published gains are measured on \textbf{data-centre}
        hardware. On a single RTX-class GPU, what actually happens?
      \end{block}
    \end{column}
    \begin{column}{0.42\textwidth}
      \scriptsize
      \begin{tabular}{@{}ll@{}}
        \toprule
        \multicolumn{2}{l}{\textbf{Open deployment questions}}\\
        \midrule
        Which draft size?      & 0.5B? 1.5B? \\
        Which proposal length? & $k=4$? $8$? $16$? \\
        Greedy or sampling?    & speed vs.\ quality \\
        Does it transfer?      & desktop $\to$ laptop \\
        \bottomrule
      \end{tabular}
    \end{column}
  \end{columns}
\end{frame}
\note{%
\textbf{[0:20--1:05 | Motivation]}\\[2pt]
Frame the bottleneck physically: every token costs a complete pass over the
target model's weights, so on a single stream you are moving gigabytes of
weights per token and the GPU's compute units are mostly idle. That is why
batch-1 interactive decoding is memory-bandwidth bound.\\[4pt]
Then the pivot: ``speculative decoding is the well-known answer, and the
literature reports two-to-five-times speedups. But look at where those numbers
come from --- A100s, H100s, B200s. The practitioner running a 3B or 4B model on
one RTX card has no reliable guidance.''\\[4pt]
Point at the four questions on the right --- these are exactly the knobs
someone has to set on day one, and say we will answer all four with
measurements.\\[4pt]
\emph{Do not} over-explain memory bandwidth; one sentence is enough.
}

% ================================================================
\begin{frame}{Speculative decoding in one slide}
  \begin{columns}[T,onlytextwidth]
    \begin{column}{0.50\textwidth}
      \textbf{Draft $\to$ verify $\to$ commit}
      \begin{enumerate}\itemsep2pt
        \item A small \textbf{draft} model proposes $k$ tokens
              autoregressively.
        \item The large \textbf{target} verifies all $k$ in
              \alert{one} parallel forward pass.
        \item Commit the accepted prefix; add one correction token on the
              first mismatch.
      \end{enumerate}
      \vspace{3pt}
      Modified rejection sampling keeps the output distribution
      \emph{identical} to the target's.
      \\[2pt]
      \tiny Leviathan et al.\ 2023; Chen et al.\ 2023
    \end{column}
    \begin{column}{0.48\textwidth}
      \begin{block}{Per-cycle latency}
        \[
          L \;=\; \frac{T_{\text{draft}} + T_{\text{verify}}}{\tau},
          \qquad
          T_{\text{draft}} = k \cdot t_{\text{draft step}}
        \]
        $\tau \in [1,\,k{+}1]$ accepted tokens per cycle.
      \end{block}
      \vspace{-2pt}
      \begin{itemize}\itemsep2pt
        \item Numerator grows \alert{linearly} in $k$.
        \item Denominator $\tau$ saturates as acceptance $\alpha$ falls.
        \item[$\Rightarrow$] There must be a $k$ beyond which it \emph{loses}.
      \end{itemize}
    \end{column}
  \end{columns}
\end{frame}
\note{%
\textbf{[1:05--2:00 | Method recap]}\\[2pt]
Keep this brisk --- most of the room knows the algorithm. Walk the three steps
once, then land hard on the equation, because it is the analytical spine of
every result that follows.\\[4pt]
Say it plainly: ``the cost in the numerator grows linearly with $k$, because
you run the draft model $k$ times. The payoff in the denominator is $\tau$, the
accepted tokens per cycle, and $\tau$ saturates --- once acceptance drops off,
extra proposals are wasted work. So the theory already tells you there is a
crossover. Our contribution is showing \emph{where} it sits on real consumer
hardware, and that it sits much lower than people assume.''\\[4pt]
Stress the correctness guarantee once --- ``this is lossless with respect to
the target distribution'' --- because it pre-empts a question, and because it
makes the quality result later more interesting.
}

% ================================================================
\begin{frame}{What is missing, and what we ask}
  \begin{columns}[T,onlytextwidth]
    \begin{column}{0.44\textwidth}
      \textbf{Three gaps in the literature}
      \begin{itemize}\itemsep3pt
        \item No cross-model \emph{and} cross-hardware validation at the
              consumer tier.
        \item Speedups reported as bare means --- no sample- or seed-level
              dispersion.
        \item Rarely any statistical test against a matched autoregressive
              baseline.
      \end{itemize}
      \vspace{4pt}
      {\small This is a \emph{characterization} paper, not a new decoding
      algorithm.}
    \end{column}
    \begin{column}{0.54\textwidth}
      \begin{block}{Research questions}
        \small
        \begin{description}\itemsep2pt
          \item[RQ1] How much end-to-end speedup on consumer RTX hardware?
          \item[RQ2] How do $k$ and acceptance \emph{jointly} set net latency?
          \item[RQ3] How sensitive is the trade-off to deterministic
                vs.\ stochastic decoding?
          \item[RQ4] Do rankings and gains transfer RTX4090 $\to$
                RTX5090-laptop?
        \end{description}
      \end{block}
    \end{column}
  \end{columns}
\end{frame}
\note{%
\textbf{[2:00--2:40 | Gap and RQs]}\\[2pt]
Be explicit and unapologetic about the contribution type: ``we are not
proposing a faster drafter. We are asking whether the existing, unmodified
draft--verify loop delivers on the hardware most people actually own, and we
are answering it with the statistical care that deployment decisions
deserve.''\\[4pt]
Read RQ1 and RQ4 aloud; gesture at RQ2 and RQ3. RQ4 is the one the audience
will remember, so give it a beat: ``and the fourth question is the one that
surprised us.''\\[4pt]
\emph{If running long:} skip the three gaps, go straight to the RQ block ---
saves about 20 seconds.
}

% ================================================================
\section{Protocol}

\begin{frame}{Experimental protocol}
  \begin{columns}[T,onlytextwidth]
    \begin{column}{0.49\textwidth}
      \small
      \textbf{Models} (same-family pairings)
      \begin{itemize}\itemsep0pt
        \item Qwen2.5-3B-Instruct $\leftarrow$ 0.5B draft
        \item Qwen3-4B-Instruct $\leftarrow$ 0.6B draft
        \item FP16, matched tokenizers, KV cache
      \end{itemize}
      \vspace{3pt}
      \textbf{Hardware}
      \begin{itemize}\itemsep0pt
        \item RTX4090 desktop, 1008\,GB/s, 450\,W
        \item RTX5090 laptop, 896\,GB/s, $\leq$150\,W
      \end{itemize}
      \vspace{3pt}
      \textbf{Data} --- 1000 samples, frozen manifests:\\
      GSM8K 300 \, $\cdot$ \, MMLU 500 \, $\cdot$ \, CNN/DM 200
    \end{column}
    \begin{column}{0.49\textwidth}
      \begin{block}{Fixed-policy grid}
        \small
        $k \in \{4, 8, 16\}$ $\times$ \{deterministic, stochastic\}\\
        $\Rightarrow$ 6 configurations \textbf{per family per device},\\
        each against a \alert{regime-matched} AR baseline.
        \\[3pt]
        Deterministic: $T{=}0$, top-$p{=}1$.\quad
        Stochastic: $T{=}0.7$, top-$p{=}0.9$.
      \end{block}
      \begin{block}{Controls}
        \small
        Seed 42 primary; repeat seeds 123/999 at $k{=}4$.
        Frozen sample-ID manifests, identical prompts and length limits,
        warm-up before measurement, serial execution.
      \end{block}
    \end{column}
  \end{columns}
\end{frame}
\note{%
\textbf{[2:40--3:40 | Protocol]}\\[2pt]
The message here is \emph{pairing}. Every speculative run is compared to a
baseline that used the same prompts, the same sample IDs, the same length
limits and the same sampling regime --- so speedup can be computed per sample,
not just as a ratio of two independent averages. That is what makes the paired
$t$-tests later legitimate.\\[4pt]
Note the deliberate choice of same-family pairings: matching tokenizers means
no vocabulary-alignment overhead contaminating the measurement. Cross-family
pairing is a different study.\\[4pt]
On hardware, flag the counter-intuitive line: the laptop 5090 is the
\emph{newer} architecture --- Blackwell versus Ada --- but it has lower memory
bandwidth and roughly a third of the power budget. Plant that seed now; it pays
off at RQ4.\\[4pt]
Be honest about the seed budget: full grid at seed 42, repeats only at $k{=}4$.
Somebody will ask.
}

% ================================================================
\section{Results}

\begin{frame}{RQ1 --- On the RTX4090, it works}
  \centering
  \small
  \begin{tabular}{l cc cc cc}
    \toprule
    & \multicolumn{2}{c}{$k=4$} & \multicolumn{2}{c}{$k=8$}
    & \multicolumn{2}{c}{$k=16$}\\
    \cmidrule(lr){2-3}\cmidrule(lr){4-5}\cmidrule(lr){6-7}
    \textbf{Speedup $S$} & det & stoch & det & stoch & det & stoch\\
    \midrule
    Qwen2.5 (3B $\leftarrow$ 0.5B) & \best{3.0674} & 2.9394 & 2.4542 & 2.4063 & 1.8050 & 1.8185\\
    Qwen3 \; (4B $\leftarrow$ 0.6B) & \best{2.8548} & 2.4798 & 1.9648 & 1.6869 & 1.1091 & \bad{0.9719}\\
    \midrule
    Qwen3 acceptance $\alpha$ & 0.336 & 0.267 & 0.238 & 0.192 & 0.129 & 0.108\\
    Qwen3 block $\Beff$       & 1.33  & 1.06  & 1.86  & 1.50  & 1.96  & 1.64\\
    \bottomrule
  \end{tabular}
  \vspace{6pt}
  \begin{columns}[T]
    \begin{column}{0.48\textwidth}
      \small
      \begin{itemize}\itemsep1pt
        \item 11 of 12 configurations beat baseline.
        \item Both families peak at \best{deterministic $k{=}4$}.
      \end{itemize}
    \end{column}
    \begin{column}{0.48\textwidth}
      \small
      \begin{itemize}\itemsep1pt
        \item Qwen3 stochastic $k{=}16$ is the one loss (\bad{$0.97\times$}).
        \item Long proposals can erase the gain entirely.
      \end{itemize}
    \end{column}
  \end{columns}
\end{frame}
\note{%
\textbf{[3:40--4:55 | RQ1 headline]}\\[2pt]
This is the good-news slide --- deliver it confidently, then immediately
complicate it.\\[4pt]
Lead with the peak: ``on the RTX4090 we get just over three times on Qwen2.5
and just under three times on Qwen3, both at deterministic $k{=}4$. Eleven of
our twelve RTX4090 configurations beat their baseline.''\\[4pt]
Then the hook into RQ2: ``but notice which column the best number is in. It is
the \emph{smallest} $k$ we tested. Every step to a longer proposal costs us
speed. And in the bottom-right cell, Qwen3 stochastic at $k{=}16$, speculative
decoding is actually \emph{slower} than just running the target model
normally.''\\[4pt]
Also point at the two bottom rows now, because they set up the next slide:
acceptance $\alpha$ falls by roughly a factor of three from $k{=}4$ to $k{=}16$,
while the effective block size only rises from 1.3 to 2.0.\\[4pt]
\emph{Pointer discipline:} touch three cells only --- the two bold peaks and
the 0.9719. Do not read the table out.
}

% ================================================================
\begin{frame}{RQ2 --- Bigger blocks, slower generation}
  \begin{columns}[T,onlytextwidth]
    \begin{column}{0.55\textwidth}
      \panelA[\linewidth]
    \end{column}
    \begin{column}{0.43\textwidth}
      \small
      \textbf{Mean speedup, $k{=}4 \to k{=}16$}
      \begin{itemize}\itemsep0pt
        \item Qwen2.5: $2.9164 \to 1.8118$
        \item Qwen3: \; $2.6673 \to 1.0405$
      \end{itemize}
      \vspace{3pt}
      \textbf{Token timing degrades too}\\
      {\scriptsize (Qwen3, RTX4090, deterministic)}
      \begin{itemize}\itemsep0pt
        \item TTFT: \; $132.6 \to 349.6$\,ms
        \item TPOT: \; $41.2 \to 94.2$\,ms/token
      \end{itemize}
      \vspace{3pt}
      \begin{block}{}
        \small $\Beff$ \emph{rises} while latency \emph{worsens}: the
        overhead inflection arrives \alert{before $k{=}8$}.
      \end{block}
    \end{column}
  \end{columns}
\end{frame}
\note{%
\textbf{[4:55--5:55 | RQ2 the inflection]}\\[2pt]
This is the analytically interesting result, so slow down here.\\[4pt]
``Our hypothesis going in --- H2 --- was the textbook one: $k{=}8$ should beat
$k{=}4$, and $k{=}16$ might start to lose. That is wrong on this hardware. The
curve is monotonically decreasing from the smallest $k$ we tested. The
inflection has already happened before $k{=}8$.''\\[4pt]
Explain \emph{why} using the equation from earlier: the effective block size
does go up --- you really are committing more tokens per verify step --- but it
goes up sub-linearly while the drafting cost goes up linearly. Acceptance is
only 34\% at $k{=}4$ and collapses to 13\% at $k{=}16$, so most of those extra
$k$ draft steps are pure waste.\\[4pt]
The token-timing numbers on the right are the practitioner's punchline: at
$k{=}16$ you have nearly tripled time-to-first-token. For an interactive
application that is the worst possible place to spend latency --- the user sees
it directly.\\[4pt]
\emph{Implication to state:} on this class of hardware, tune $k$ downward, and
$k{=}4$ may not even be the floor.
}

% ================================================================
\begin{frame}{RQ3 --- Deterministic decoding is the safer default}
  \begin{columns}[T,onlytextwidth]
    \begin{column}{0.52\textwidth}
      \scriptsize
      \begin{tabular}{@{}lccc@{}}
        \toprule
        \textbf{Qwen3, RTX4090} & $k{=}4$ & $k{=}8$ & $k{=}16$\\
        \midrule
        Deterministic $S$ & 2.8548 & 1.9648 & 1.1091\\
        Stochastic $S$    & 2.4798 & 1.6869 & 0.9719\\
        \midrule
        Advantage         & \best{$+0.3750$} & $+0.2779$ & $+0.1372$\\
        \bottomrule
      \end{tabular}
      \vspace{8pt}
      \small
      \begin{itemize}\itemsep3pt
        \item Greedy drafting agrees with greedy verification more often
              $\Rightarrow$ higher $\alpha$ (0.336 vs.\ 0.267 at $k{=}4$).
        \item The advantage \emph{shrinks} with $k$ --- once acceptance is
              low, regime hardly matters.
      \end{itemize}
    \end{column}
    \begin{column}{0.46\textwidth}
      \begin{block}{Not universal}
        \small
        Qwen2.5 follows the same ordering at $k{=}4$ and $k{=}8$, but reverses
        slightly at $k{=}16$ ($1.8050$ det vs.\ $1.8185$ stoch).
      \end{block}
      \vspace{2pt}
      {\small\textbf{Takeaway.} If you need sampling, expect to pay
      $\sim$13\% of your speedup for it at low $k$.}
    \end{column}
  \end{columns}
\end{frame}
\note{%
\textbf{[5:55--6:30 | RQ3 regime]}\\[2pt]
Short slide --- this is the least surprising result, so move through it.\\[4pt]
The mechanism is the thing worth saying out loud: ``when both draft and target
decode greedily, the draft's proposal is compared against a deterministic
target continuation, so agreement is high. Turn sampling on for both and you
introduce two independent sources of randomness, acceptance drops from about
34\% to 27\%, and you lose roughly 13\% of your speedup.''\\[4pt]
Be careful to state the caveat rather than hiding it: at $k{=}16$ Qwen2.5
actually reverses, marginally. It is a small effect but we report it --- the
preference is a strong tendency, not a law.\\[4pt]
\emph{If running long:} this is the first slide to cut. State the one-line
takeaway --- ``deterministic wins at matched $k$, by a margin that shrinks as
$k$ grows'' --- and click through.
}

% ================================================================
\begin{frame}{RQ4 --- The same policy, a different GPU, and it all inverts}
  \begin{columns}[T,onlytextwidth]
    \begin{column}{0.53\textwidth}
      \panelB[\linewidth]
    \end{column}
    \begin{column}{0.45\textwidth}
      \small
      \textbf{Qwen3 on RTX5090-laptop}
      \scriptsize
      \begin{tabular}{@{}lcc@{}}
        \toprule
        Config & $S$ & $\alpha$\\
        \midrule
        $k{=}4$, det   & \bad{0.6743} & 0.335\\
        $k{=}4$, stoch & \bad{0.5924} & 0.285\\
        $k{=}8$, det   & \bad{0.4613} & 0.254\\
        $k{=}8$, stoch & \bad{0.3995} & 0.227\\
        $k{=}16$, det  & \bad{0.2588} & 0.180\\
        $k{=}16$, stoch& \bad{0.2222} & 0.154\\
        \bottomrule
      \end{tabular}
      \vspace{4pt}
      \small
      \begin{itemize}\itemsep1pt
        \item \alert{All six are slower than baseline.}
        \item Acceptance \emph{equals or beats} RTX4090 at every $k$ ---
              the loss is in execution cost, not the algorithm.
        \item Ordering survives: low $k$, deterministic.
      \end{itemize}
    \end{column}
  \end{columns}
\end{frame}
\note{%
\textbf{[6:30--7:50 | RQ4 portability --- the core result]}\\[2pt]
This is the slide the talk exists for. Pause before starting it.\\[4pt]
``We took the identical code, the identical models, the identical frozen
manifests, and moved to an RTX5090 laptop --- a \emph{newer} architecture. Every
single configuration is now slower than simply running the target model
autoregressively. The best case, deterministic $k{=}4$, gives us 0.67 --- we have
made inference one and a half times slower.''\\[4pt]
Then the diagnostic point, which is what makes this a finding rather than an
anecdote: ``acceptance is not the problem. At $k{=}4$ it is identical --- 0.336
against 0.335. At $k{=}8$ and $k{=}16$ the laptop actually accepts \emph{more}
--- 0.25 against 0.24, and 0.18 against 0.13. So the algorithm is doing at
least as well as on the desktop, and it is still slower everywhere. What
changed is the runtime economics: on this stack the draft model's per-step cost
relative to the target's verification cost is unfavourable, and speculative
decoding needs that ratio to be small.''\\[4pt]
Land the deployment rule: ``speedup is a property of the complete
model--hardware--runtime stack. It is not a property of the algorithm, and it is
certainly not a property of the GPU generation --- the newer card lost.''\\[4pt]
Finally, the constructive half: the \emph{ranking} did transfer. Low $k$ beats
high $k$, deterministic beats stochastic, on both devices. So policy
exploration on one machine is transferable; absolute claims are not.
}
\note{%
\textbf{[RQ4, continued]}\\[2pt]
\emph{Anticipate the question:} we do not claim to have isolated the cause. Our
candidate explanation is host-side orchestration and graph-capture behaviour ---
see the CUDA-graph note in the paper --- but we did not measure it, so we report
it as context, not as a cause. Say that before someone asks.\\[4pt]
\emph{Second likely question:} ``is the laptop card simply weaker?'' Partly ---
lower bandwidth and roughly a third of the power budget. But $S$ is a
\emph{ratio} against that same device's own autoregressive baseline, so a
uniformly slower card should not by itself push the ratio below one. The
backup hardware slide has the numbers.
}

% ================================================================
\begin{frame}{Is the effect real? Statistics and stability}
  \begin{columns}[T,onlytextwidth]
    \begin{column}{0.52\textwidth}
      \small
      \textbf{Paired one-sided $t$-tests}, $H_0{:}\ S \le 1$,\\
      paired by sample, Holm--Bonferroni across 6 configs
      (Qwen3, RTX4090, $n{=}1000$):
      \vspace{2pt}
      \scriptsize
      \begin{tabular}{@{}lcc@{}}
        \toprule
        Config & $d_S$ & $p_S$\\
        \midrule
        det $k{=}4,8,16$   & 1.18 / 1.10 / 0.29 & all $<10^{-18}$\\
        stoch $k{=}4,8$    & 1.24 / 1.13       & both $<10^{-179}$\\
        stoch $k{=}16$     & \bad{$-0.08$}     & \bad{$0.996$} \\
        \bottomrule
      \end{tabular}
      \vspace{6pt}
      \small
      \textbf{Dispersion is large.} Per-sample $S$ (det):
      $2.7561\pm0.5950$ at $k{=}4$ $\to$ $1.3549\pm0.6533$ at $k{=}16$.
      \\[2pt]
      {\scriptsize At $k{=}16$ the standard deviation is half the mean ---
      prompt-to-prompt variation dominates.}
    \end{column}
    \begin{column}{0.46\textwidth}
      \panelD[\linewidth]
      \vspace{-2pt}
      {\scriptsize Seed 999 vs.\ 123 at $k{=}4$: deterministic within
      $1.5\%$; stochastic $0.4\%$--$5.0\%$.}
    \end{column}
  \end{columns}
\end{frame}
\note{%
\textbf{[7:50--8:40 | Evidence quality]}\\[2pt]
Purpose of this slide: pre-empt the methodological objections, and show we are
not over-claiming.\\[4pt]
``Five of six RTX4090 configurations are significant after Holm--Bonferroni
correction, with large effect sizes. The sixth --- stochastic $k{=}16$ --- is
not significant, and its effect size is slightly negative, which is exactly
consistent with its sub-baseline mean. We report it as a null, not as noise.''\\[4pt]
The dispersion point matters more than the $p$-values: at $k{=}16$ the standard
deviation across prompts is roughly half the mean. A single reported average
speedup hides an enormous spread --- your slowest prompts may see no benefit at
all. This is one of the things the literature routinely omits.\\[4pt]
On the right panel: repeat seeds at $k{=}4$. Deterministic runs move by at most
1.5\% between seeds; stochastic by up to 5\%. Regime ordering is preserved
everywhere.\\[4pt]
\emph{Say the limitation before the reviewer does:} we only have repeat seeds at
$k{=}4$. We do not claim seed stability at $k{=}8$ or $k{=}16$.
}

% ================================================================
\begin{frame}{Speed is not free: quality on CNN/DailyMail}
  \centering
  \small
  \begin{tabular}{lccc}
    \toprule
    \textbf{Deterministic, $n{=}200$} & ROUGE-L & BLEU & BERTScore F1\\
    \midrule
    Autoregressive baseline & 19.58 & 3.20 & 0.793\\
    Speculative, $k{=}4$    & \bad{11.50} & \bad{1.39} & \bad{0.710}\\
    Speculative, $k{=}16$   & 14.92 & 2.04 & 0.741\\
    \bottomrule
  \end{tabular}
  \vspace{8pt}
  \begin{columns}[T]
    \begin{column}{0.47\textwidth}
      \small
      \begin{itemize}\itemsep2pt
        \item All three metrics agree: a real quality cost on the
              summarization task.
        \item It \alert{partially recovers} as $k$ grows --- moving
              \emph{against} the speed trend.
      \end{itemize}
    \end{column}
    \begin{column}{0.47\textwidth}
      \begin{block}{}
        \small The fastest configuration is also the weakest on
        summarization. There is a genuine speed--quality frontier here, not
        free acceleration.
      \end{block}
    \end{column}
  \end{columns}
\end{frame}
\note{%
\textbf{[8:40--9:25 | Quality trade-off]}\\[2pt]
Deliver this as an honest complication of our own headline number.\\[4pt]
``Earlier I said the peak was deterministic $k{=}4$. On summarization, that is
also our \emph{worst} quality point --- ROUGE-L drops from 19.6 to 11.5, and BLEU
and BERTScore move the same way. Three independent metrics agreeing means this
is not a metric artefact.''\\[4pt]
The direction is the interesting part: quality \emph{recovers} as $k$ increases,
which is the opposite direction to speed. So the two axes genuinely trade off
against each other across our grid, and the operating point is a choice, not a
default.\\[4pt]
\emph{Expected question --- be ready:} ``isn't speculative decoding supposed to
be distribution-preserving?'' Answer: exactly, and that is why this is worth
reporting. The guarantee is about the sampling distribution under the idealized
scheme; what we measure is an end-to-end implementation on a bounded generation
budget, where accepted-prefix commitment and length limits interact. We flag it
as an observed effect in our implementation and do not claim it refutes the
theory. Do not oversell this --- it is a limitation we are surfacing.\\[4pt]
Note that GSM8K and MMLU are accuracy tasks and are covered by our success
criteria; this slide is the summarization result specifically.
}

% ================================================================
\section{Wrap-up}

\begin{frame}{What we can and cannot claim}
  \begin{columns}[T,onlytextwidth]
    \begin{column}{0.48\textwidth}
      \textbf{Hypotheses}
      \small
      \begin{itemize}\itemsep2pt
        \item[\textbf{H1}] $S>1$ --- \best{supported} on RTX4090
              (11/12), \bad{rejected} on RTX5090-laptop.
        \item[\textbf{H2}] $k{=}8 > k{=}4$ --- \bad{rejected}; speedup
              falls monotonically.
        \item[\textbf{H3}] $\Beff$ up, latency worse --- \best{supported}.
        \item[\textbf{H4}] Deterministic faster --- \best{mostly supported}
              (one reversal).
      \end{itemize}
    \end{column}
    \begin{column}{0.48\textwidth}
      \textbf{Threats to validity}
      \small
      \begin{itemize}\itemsep2pt
        \item One PyTorch/Transformers stack; results do not extrapolate to
              other runtimes or GPUs.
        \item Same-family Qwen pairings only.
        \item No long-context, safety, or factual-consistency evaluation.
        \item Repeat seeds cover $k{=}4$ only.
        \item CUDA-graph capture differs across run families --- reported as
              context, \emph{not} as a measured cause.
      \end{itemize}
    \end{column}
  \end{columns}
\end{frame}
\note{%
\textbf{[9:25--9:55 | Scorecard and limits]}\\[2pt]
Move quickly --- this slide is mostly for the record and for the reviewers in
the room. Do not read every bullet.\\[4pt]
Highlight two things only. First, H2: ``we predicted $k{=}8$ would beat $k{=}4$
and it does not --- the overhead inflection is earlier than the folklore
suggests.'' Second, the limitations column exists because we want the claim
bounded: this is one runtime stack, one model family, one benchmark suite.\\[4pt]
The CUDA-graph bullet is the intellectually honest one --- we saw a difference in
capture behaviour between run families, and it is a plausible contributor to the
RTX5090 result, but we did not isolate it experimentally, so we refuse to call
it the cause.\\[4pt]
\emph{Timing checkpoint:} you should be at roughly 9:30 here. If you are past
10:30, skip straight to the takeaways.
}

% ================================================================
\begin{frame}[standout]
  \Large
  Policy ranking is portable. \\[6pt]
  Absolute speedup is not.
\end{frame}
\note{%
\textbf{[9:55--10:05 | The one line]}\\[2pt]
Pause. Let it sit for two seconds. This is the sentence you want people to
repeat to a colleague afterwards. Do not add anything to it --- click through
to the takeaways.
}

% ================================================================
\begin{frame}{Takeaways}
  \begin{enumerate}\itemsep6pt
    \item \textbf{Use a short proposal.} On both consumer GPUs and both model
          families, $k{=}4$ dominates $k{=}8$ and $k{=}16$ --- the overhead
          inflection arrives earlier than expected.
    \item \textbf{Prefer deterministic decoding} at matched $k$; sampling
          costs roughly 13\% of the speedup at $k{=}4$.
    \item \textbf{Benchmark on your target stack.} An RTX4090 gave up to
          $3.07\times$; the same code on an RTX5090 laptop was
          \emph{sub-baseline everywhere}, with unchanged acceptance.
    \item \textbf{Report dispersion and quality, not just a mean speedup.}
          Per-sample spread is large, and summarization quality moves against
          the speed trend.
  \end{enumerate}
  \vspace{6pt}
  \centering
  \small Frozen manifests, matched baselines, and raw per-sample artifacts
  support reproduction.
\end{frame}
\note{%
\textbf{[10:05--10:40 | Close]}\\[2pt]
Four sentences, one per point, then stop. Do not re-explain --- the audience has
seen all the evidence.\\[4pt]
Rule one: short proposals. Rule two: greedy if you can afford it. Rule three,
and this is the one to emphasize with your voice: measure on the machine you
will actually deploy on, because we have a concrete counter-example where a
newer GPU turned a three-times win into a one-and-a-half-times loss. Rule four:
report the spread and the quality, because the mean speedup on its own is not a
deployment decision.\\[4pt]
Close on future work in one breath: full grid across seeds, broader hardware and
model coverage, longer contexts and factuality-sensitive metrics.\\[4pt]
``Thank you --- I'm happy to take questions.''
}

% ================================================================
\appendix

\begin{frame}[standout]
  Backup slides
\end{frame}
\note{Backup material follows. Jump here only if asked.}

\begin{frame}{Backup --- full Qwen3 RTX4090 grid}
  \centering
  \small
  \begin{tabular}{lcccccc}
    \toprule
    \textbf{Config} & $S$ & $\alpha$ & $\Beff$ & $T_{\text{mean}}$ (s)
    & TTFT (ms) & TPOT (ms)\\
    \midrule
    $k{=}4$, det   & 2.8548 & 0.3362 & 1.33 & 4.8881 & 132.57 & 41.23\\
    $k{=}4$, stoch & 2.4798 & 0.2672 & 1.06 & 5.6196 & 134.54 & 45.51\\
    $k{=}8$, det   & 1.9648 & 0.2378 & 1.86 & 7.1024 & 204.99 & 58.76\\
    $k{=}8$, stoch & 1.6869 & 0.1916 & 1.50 & 8.2612 & 207.82 & 64.67\\
    $k{=}16$, det  & 1.1091 & 0.1292 & 1.96 & 12.5817 & 349.58 & 94.24\\
    $k{=}16$, stoch& 0.9719 & 0.1082 & 1.64 & 14.3390 & 354.82 & 106.03\\
    \bottomrule
  \end{tabular}
  \vspace{6pt}

  \small Target: Qwen3-4B-Instruct; draft: Qwen3-0.6B-Instruct; $n{=}1000$.
\end{frame}
\note{Use if asked for the full Qwen3 numbers, or for TTFT/TPOT detail.}

\begin{frame}{Backup --- full Qwen2.5 RTX4090 grid}
  \centering
  \small
  \begin{tabular}{lcccccc}
    \toprule
    \textbf{Config} & $S$ & $\alpha$ & $\Beff$ & $T_{\text{mean}}$ (s)
    & TTFT (ms) & TPOT (ms)\\
    \midrule
    $k{=}4$, det   & 3.0674 & 0.4212 & 1.66 & 3.7210 & 102.35 & 29.06\\
    $k{=}4$, stoch & 2.9394 & 0.4080 & 1.61 & 3.9631 & 105.26 & 30.89\\
    $k{=}8$, det   & 2.4542 & 0.3515 & 2.72 & 4.6508 & 162.76 & 36.73\\
    $k{=}8$, stoch & 2.4063 & 0.3444 & 2.67 & 4.8411 & 165.95 & 38.04\\
    $k{=}16$, det  & 1.8050 & 0.2539 & 3.76 & 6.3234 & 277.43 & 49.34\\
    $k{=}16$, stoch& 1.8185 & 0.2516 & 3.73 & 6.4059 & 278.85 & 49.81\\
    \bottomrule
  \end{tabular}
  \vspace{6pt}

  \small Target: Qwen2.5-3B-Instruct; draft: Qwen2.5-0.5B-Instruct; $n{=}1000$.\\
  Higher $\alpha$ and $\Beff$ than Qwen3 at every $k$.
\end{frame}
\note{%
Qwen2.5 accepts more than Qwen3 at every $k$ --- $\alpha=0.42$ vs.\ $0.34$ at
$k{=}4$, and effective block size 3.76 vs.\ 1.96 at $k{=}16$. That is why
Qwen2.5 never falls below baseline on the RTX4090 while Qwen3 does. Likely a
smaller capability gap between the 3B target and 0.5B draft than between the
4B target and 0.6B draft.
}

\begin{frame}{Backup --- acceptance portability}
  \centering
  \panelC[0.62\linewidth]
  \vspace{2pt}

  \small Qwen3 acceptance on the laptop is identical at $k{=}4$ (0.335 vs.\
  0.336) and \emph{higher} at $k{=}8$ and $k{=}16$ (0.254 vs.\ 0.238; 0.180
  vs.\ 0.129). The algorithm behaves at least as well; the
  \alert{runtime economics} differ.
\end{frame}
\note{%
The single most useful backup slide for the RQ4 discussion. If someone
challenges the RTX5090 result --- ``did you misconfigure it?'' --- this is the
answer: acceptance is a model-level quantity; it is unchanged at $k{=}4$ and
actually \emph{better} at higher $k$, so the speculative loop is doing at least
as well as it did on the desktop. The loss is in per-step execution cost, not
in the algorithm.
}

\begin{frame}{Backup --- success criteria and metrics}
  \begin{columns}[T,onlytextwidth]
    \begin{column}{0.48\textwidth}
      \small
      \textbf{Stated success criteria}
      \begin{itemize}\itemsep2pt
        \item Primary: mean $S \ge 1.3\times$ with Holm--Bonferroni-corrected
              $p_S < 0.05$.
        \item Primary: absolute quality drop $\le 1.0$ point on GSM8K and
              MMLU, deterministic regime.
        \item Secondary: best $(k, \text{regime})$ under quality constraints;
              dispersion; seed stability; cross-hardware ranking transfer.
      \end{itemize}
    \end{column}
    \begin{column}{0.48\textwidth}
      \small
      \textbf{Key metric definitions}
      \begin{itemize}\itemsep2pt
        \item $S$ = baseline latency / speculative latency
        \item $\alpha$ = accepted / proposed draft tokens
        \item $\Beff$ = accepted draft tokens / verify steps
        \item TTFT = first token time $-$ request start
        \item TPOT = (completion $-$ first token) / output tokens
        \item $d_S$ = Cohen's $d$, paired baseline vs.\ speculative latency
      \end{itemize}
    \end{column}
  \end{columns}
\end{frame}
\note{Reference slide for methodology questions.}

\begin{frame}{Backup --- hardware profiles}
  \centering
  \small
  \begin{tabular}{lll}
    \toprule
    \textbf{Component} & \textbf{RTX4090 desktop} & \textbf{RTX5090 laptop}\\
    \midrule
    GPU memory        & 24\,GB GDDR6X & 24\,GB GDDR7\\
    Memory bandwidth  & \best{1008\,GB/s} & 896\,GB/s\\
    CUDA cores        & \best{16384} & 10496\\
    Architecture      & Ada Lovelace & \best{Blackwell}\\
    Power profile     & \best{450\,W TGP class} & up to 150\,W TGP (OEM-dependent)\\
    Software stack    & \multicolumn{2}{c}{PyTorch + Transformers, pinned per run family}\\
    \bottomrule
  \end{tabular}
  \vspace{8pt}

  \small The laptop part is the newer architecture but has lower bandwidth and
  roughly \alert{one third of the power budget}.
\end{frame}
\note{%
The obvious question after RQ4 is ``is the 5090 laptop just a weaker card?''
Partly --- lower bandwidth, far lower sustained power. But note that this affects
the autoregressive baseline too, and $S$ is a \emph{ratio} against that same
device's baseline. So a uniformly slower card should not by itself flip the
ratio below 1. That is why we point at the draft/target cost ratio and the
runtime path rather than at raw card performance.
}

\end{document}
"
%    Dual screen :  set \def\NOTESMODE{} -> change option below
% ===============================================================
\documentclass[aspectratio=169, 11pt]{beamer}

% ---------------------------------------------------------------
% Shared preamble for the ICECCME 2026 presentation
% ---------------------------------------------------------------
\usetheme[progressbar=frametitle, numbering=fraction, block=fill]{metropolis}
\usepackage[T1]{fontenc}
\usepackage[utf8]{inputenc}
\usepackage{amsmath,amssymb}
\usepackage{booktabs}
\usepackage{graphicx}
\usepackage{xcolor}
\usepackage{appendixnumberbeamer}

% Colour-blind-safe palette, matching the paper figures.
\definecolor{cbBlue}{HTML}{0072B2}   % Qwen2.5 / RTX4090
\definecolor{cbOrange}{HTML}{D55E00} % Qwen3 / stochastic
\definecolor{cbGreen}{HTML}{009E73}  % RTX5090-laptop
\definecolor{cbGrey}{HTML}{555555}

\setbeamercolor{progress bar}{fg=cbBlue}
\setbeamercolor{alerted text}{fg=cbOrange}
\setbeamerfont{note page}{size=\small}

% Compact tables that survive a projector.
\setlength{\tabcolsep}{4.2pt}
\renewcommand{\arraystretch}{1.08}

% Shorthands used throughout the deck.
\newcommand{\up}{\textcolor{cbBlue}{$\uparrow$}}
\newcommand{\dn}{\textcolor{cbOrange}{$\downarrow$}}
\newcommand{\best}[1]{\textcolor{cbBlue}{\textbf{#1}}}
\newcommand{\bad}[1]{\textcolor{cbOrange}{\textbf{#1}}}
\newcommand{\Beff}{B_{\text{eff}}}

% Panel clipping for figures/specdec_configs_vs_baseline_metrics.pdf
% (a 2x2 matplotlib grid, 545.9 x 393.2 pt).
\newcommand{\panelA}[1][0.92\linewidth]{%
  \includegraphics[width=#1, trim=9 202 271 5, clip]%
  {../figures/specdec_configs_vs_baseline_metrics.pdf}}
\newcommand{\panelB}[1][0.92\linewidth]{%
  \includegraphics[width=#1, trim=271 202 22 5, clip]%
  {../figures/specdec_configs_vs_baseline_metrics.pdf}}
\newcommand{\panelC}[1][0.92\linewidth]{%
  \includegraphics[width=#1, trim=5 8 267 199, clip]%
  {../figures/specdec_configs_vs_baseline_metrics.pdf}}
\newcommand{\panelD}[1][0.92\linewidth]{%
  \includegraphics[width=#1, trim=271 8 5 199, clip]%
  {../figures/specdec_configs_vs_baseline_metrics.pdf}}


% --- notes switch -----------------------------------------------
\ifdefined\NOTESMODE
  \setbeameroption{show only notes}
  % Notes-only build: drop the slide thumbnail and shrink the type so that
  % even the longest note fits on a single page.
  \setbeamertemplate{note page}[plain]
  \setbeamerfont{note page}{size=\footnotesize}
\else
  \setbeameroption{hide notes}
  % For a presenter-display setup, comment the line above and use:
  % \setbeameroption{show notes on second screen=right}
\fi
% ----------------------------------------------------------------

\title{Characterizing Speculative Decoding Dynamics for Large
       Language Models on Consumer-Class GPUs}
\author{Samsu Sempena \and Zhuohan Cui}
\institute{The University of Sydney}
\date{ICECCME 2026 \\ 15--17 October 2026, Bali, Indonesia}

\begin{document}

% ================================================================
\begin{frame}[plain]
  \titlepage
\end{frame}
\note{%
\textbf{[0:00--0:20 | Title]}\\[2pt]
``Good morning. I'm Samsu Sempena from the University of Sydney, and this is
joint work with Zhuohan Cui. Our paper characterizes what speculative decoding
actually does on consumer-class GPUs --- the kind of hardware a practitioner
has under their desk, not in a data centre.''\\[4pt]
\emph{Cue:} say the one-line thesis up front so the audience has a hook ---
``the short version is that the \emph{policy ranking} travels between GPUs,
but the \emph{speedup} does not.''\\[4pt]
\emph{Housekeeping:} confirm the clicker works, and check whether the chair
wants questions during or after.
}

% ================================================================
\section{Motivation}

\begin{frame}{Decoding is sequential, and that is the bottleneck}
  \begin{columns}[T,onlytextwidth]
    \begin{column}{0.56\textwidth}
      \begin{itemize}
        \item Autoregressive decoding needs \alert{one full forward pass of the
              target model per token}.
        \item Single-stream, interactive latency is therefore bound by
              \emph{memory traffic}, not by arithmetic.
        \item Speculative decoding is the standard fix --- and it is reported
              to be a large, free win.
      \end{itemize}
      \vspace{4pt}
      \begin{block}{The practitioner's problem}
        Almost all published gains are measured on \textbf{data-centre}
        hardware. On a single RTX-class GPU, what actually happens?
      \end{block}
    \end{column}
    \begin{column}{0.42\textwidth}
      \scriptsize
      \begin{tabular}{@{}ll@{}}
        \toprule
        \multicolumn{2}{l}{\textbf{Open deployment questions}}\\
        \midrule
        Which draft size?      & 0.5B? 1.5B? \\
        Which proposal length? & $k=4$? $8$? $16$? \\
        Greedy or sampling?    & speed vs.\ quality \\
        Does it transfer?      & desktop $\to$ laptop \\
        \bottomrule
      \end{tabular}
    \end{column}
  \end{columns}
\end{frame}
\note{%
\textbf{[0:20--1:05 | Motivation]}\\[2pt]
Frame the bottleneck physically: every token costs a complete pass over the
target model's weights, so on a single stream you are moving gigabytes of
weights per token and the GPU's compute units are mostly idle. That is why
batch-1 interactive decoding is memory-bandwidth bound.\\[4pt]
Then the pivot: ``speculative decoding is the well-known answer, and the
literature reports two-to-five-times speedups. But look at where those numbers
come from --- A100s, H100s, B200s. The practitioner running a 3B or 4B model on
one RTX card has no reliable guidance.''\\[4pt]
Point at the four questions on the right --- these are exactly the knobs
someone has to set on day one, and say we will answer all four with
measurements.\\[4pt]
\emph{Do not} over-explain memory bandwidth; one sentence is enough.
}

% ================================================================
\begin{frame}{Speculative decoding in one slide}
  \begin{columns}[T,onlytextwidth]
    \begin{column}{0.50\textwidth}
      \textbf{Draft $\to$ verify $\to$ commit}
      \begin{enumerate}\itemsep2pt
        \item A small \textbf{draft} model proposes $k$ tokens
              autoregressively.
        \item The large \textbf{target} verifies all $k$ in
              \alert{one} parallel forward pass.
        \item Commit the accepted prefix; add one correction token on the
              first mismatch.
      \end{enumerate}
      \vspace{3pt}
      Modified rejection sampling keeps the output distribution
      \emph{identical} to the target's.
      \\[2pt]
      \tiny Leviathan et al.\ 2023; Chen et al.\ 2023
    \end{column}
    \begin{column}{0.48\textwidth}
      \begin{block}{Per-cycle latency}
        \[
          L \;=\; \frac{T_{\text{draft}} + T_{\text{verify}}}{\tau},
          \qquad
          T_{\text{draft}} = k \cdot t_{\text{draft step}}
        \]
        $\tau \in [1,\,k{+}1]$ accepted tokens per cycle.
      \end{block}
      \vspace{-2pt}
      \begin{itemize}\itemsep2pt
        \item Numerator grows \alert{linearly} in $k$.
        \item Denominator $\tau$ saturates as acceptance $\alpha$ falls.
        \item[$\Rightarrow$] There must be a $k$ beyond which it \emph{loses}.
      \end{itemize}
    \end{column}
  \end{columns}
\end{frame}
\note{%
\textbf{[1:05--2:00 | Method recap]}\\[2pt]
Keep this brisk --- most of the room knows the algorithm. Walk the three steps
once, then land hard on the equation, because it is the analytical spine of
every result that follows.\\[4pt]
Say it plainly: ``the cost in the numerator grows linearly with $k$, because
you run the draft model $k$ times. The payoff in the denominator is $\tau$, the
accepted tokens per cycle, and $\tau$ saturates --- once acceptance drops off,
extra proposals are wasted work. So the theory already tells you there is a
crossover. Our contribution is showing \emph{where} it sits on real consumer
hardware, and that it sits much lower than people assume.''\\[4pt]
Stress the correctness guarantee once --- ``this is lossless with respect to
the target distribution'' --- because it pre-empts a question, and because it
makes the quality result later more interesting.
}

% ================================================================
\begin{frame}{What is missing, and what we ask}
  \begin{columns}[T,onlytextwidth]
    \begin{column}{0.44\textwidth}
      \textbf{Three gaps in the literature}
      \begin{itemize}\itemsep3pt
        \item No cross-model \emph{and} cross-hardware validation at the
              consumer tier.
        \item Speedups reported as bare means --- no sample- or seed-level
              dispersion.
        \item Rarely any statistical test against a matched autoregressive
              baseline.
      \end{itemize}
      \vspace{4pt}
      {\small This is a \emph{characterization} paper, not a new decoding
      algorithm.}
    \end{column}
    \begin{column}{0.54\textwidth}
      \begin{block}{Research questions}
        \small
        \begin{description}\itemsep2pt
          \item[RQ1] How much end-to-end speedup on consumer RTX hardware?
          \item[RQ2] How do $k$ and acceptance \emph{jointly} set net latency?
          \item[RQ3] How sensitive is the trade-off to deterministic
                vs.\ stochastic decoding?
          \item[RQ4] Do rankings and gains transfer RTX4090 $\to$
                RTX5090-laptop?
        \end{description}
      \end{block}
    \end{column}
  \end{columns}
\end{frame}
\note{%
\textbf{[2:00--2:40 | Gap and RQs]}\\[2pt]
Be explicit and unapologetic about the contribution type: ``we are not
proposing a faster drafter. We are asking whether the existing, unmodified
draft--verify loop delivers on the hardware most people actually own, and we
are answering it with the statistical care that deployment decisions
deserve.''\\[4pt]
Read RQ1 and RQ4 aloud; gesture at RQ2 and RQ3. RQ4 is the one the audience
will remember, so give it a beat: ``and the fourth question is the one that
surprised us.''\\[4pt]
\emph{If running long:} skip the three gaps, go straight to the RQ block ---
saves about 20 seconds.
}

% ================================================================
\section{Protocol}

\begin{frame}{Experimental protocol}
  \begin{columns}[T,onlytextwidth]
    \begin{column}{0.49\textwidth}
      \small
      \textbf{Models} (same-family pairings)
      \begin{itemize}\itemsep0pt
        \item Qwen2.5-3B-Instruct $\leftarrow$ 0.5B draft
        \item Qwen3-4B-Instruct $\leftarrow$ 0.6B draft
        \item FP16, matched tokenizers, KV cache
      \end{itemize}
      \vspace{3pt}
      \textbf{Hardware}
      \begin{itemize}\itemsep0pt
        \item RTX4090 desktop, 1008\,GB/s, 450\,W
        \item RTX5090 laptop, 896\,GB/s, $\leq$150\,W
      \end{itemize}
      \vspace{3pt}
      \textbf{Data} --- 1000 samples, frozen manifests:\\
      GSM8K 300 \, $\cdot$ \, MMLU 500 \, $\cdot$ \, CNN/DM 200
    \end{column}
    \begin{column}{0.49\textwidth}
      \begin{block}{Fixed-policy grid}
        \small
        $k \in \{4, 8, 16\}$ $\times$ \{deterministic, stochastic\}\\
        $\Rightarrow$ 6 configurations \textbf{per family per device},\\
        each against a \alert{regime-matched} AR baseline.
        \\[3pt]
        Deterministic: $T{=}0$, top-$p{=}1$.\quad
        Stochastic: $T{=}0.7$, top-$p{=}0.9$.
      \end{block}
      \begin{block}{Controls}
        \small
        Seed 42 primary; repeat seeds 123/999 at $k{=}4$.
        Frozen sample-ID manifests, identical prompts and length limits,
        warm-up before measurement, serial execution.
      \end{block}
    \end{column}
  \end{columns}
\end{frame}
\note{%
\textbf{[2:40--3:40 | Protocol]}\\[2pt]
The message here is \emph{pairing}. Every speculative run is compared to a
baseline that used the same prompts, the same sample IDs, the same length
limits and the same sampling regime --- so speedup can be computed per sample,
not just as a ratio of two independent averages. That is what makes the paired
$t$-tests later legitimate.\\[4pt]
Note the deliberate choice of same-family pairings: matching tokenizers means
no vocabulary-alignment overhead contaminating the measurement. Cross-family
pairing is a different study.\\[4pt]
On hardware, flag the counter-intuitive line: the laptop 5090 is the
\emph{newer} architecture --- Blackwell versus Ada --- but it has lower memory
bandwidth and roughly a third of the power budget. Plant that seed now; it pays
off at RQ4.\\[4pt]
Be honest about the seed budget: full grid at seed 42, repeats only at $k{=}4$.
Somebody will ask.
}

% ================================================================
\section{Results}

\begin{frame}{RQ1 --- On the RTX4090, it works}
  \centering
  \small
  \begin{tabular}{l cc cc cc}
    \toprule
    & \multicolumn{2}{c}{$k=4$} & \multicolumn{2}{c}{$k=8$}
    & \multicolumn{2}{c}{$k=16$}\\
    \cmidrule(lr){2-3}\cmidrule(lr){4-5}\cmidrule(lr){6-7}
    \textbf{Speedup $S$} & det & stoch & det & stoch & det & stoch\\
    \midrule
    Qwen2.5 (3B $\leftarrow$ 0.5B) & \best{3.0674} & 2.9394 & 2.4542 & 2.4063 & 1.8050 & 1.8185\\
    Qwen3 \; (4B $\leftarrow$ 0.6B) & \best{2.8548} & 2.4798 & 1.9648 & 1.6869 & 1.1091 & \bad{0.9719}\\
    \midrule
    Qwen3 acceptance $\alpha$ & 0.336 & 0.267 & 0.238 & 0.192 & 0.129 & 0.108\\
    Qwen3 block $\Beff$       & 1.33  & 1.06  & 1.86  & 1.50  & 1.96  & 1.64\\
    \bottomrule
  \end{tabular}
  \vspace{6pt}
  \begin{columns}[T]
    \begin{column}{0.48\textwidth}
      \small
      \begin{itemize}\itemsep1pt
        \item 11 of 12 configurations beat baseline.
        \item Both families peak at \best{deterministic $k{=}4$}.
      \end{itemize}
    \end{column}
    \begin{column}{0.48\textwidth}
      \small
      \begin{itemize}\itemsep1pt
        \item Qwen3 stochastic $k{=}16$ is the one loss (\bad{$0.97\times$}).
        \item Long proposals can erase the gain entirely.
      \end{itemize}
    \end{column}
  \end{columns}
\end{frame}
\note{%
\textbf{[3:40--4:55 | RQ1 headline]}\\[2pt]
This is the good-news slide --- deliver it confidently, then immediately
complicate it.\\[4pt]
Lead with the peak: ``on the RTX4090 we get just over three times on Qwen2.5
and just under three times on Qwen3, both at deterministic $k{=}4$. Eleven of
our twelve RTX4090 configurations beat their baseline.''\\[4pt]
Then the hook into RQ2: ``but notice which column the best number is in. It is
the \emph{smallest} $k$ we tested. Every step to a longer proposal costs us
speed. And in the bottom-right cell, Qwen3 stochastic at $k{=}16$, speculative
decoding is actually \emph{slower} than just running the target model
normally.''\\[4pt]
Also point at the two bottom rows now, because they set up the next slide:
acceptance $\alpha$ falls by roughly a factor of three from $k{=}4$ to $k{=}16$,
while the effective block size only rises from 1.3 to 2.0.\\[4pt]
\emph{Pointer discipline:} touch three cells only --- the two bold peaks and
the 0.9719. Do not read the table out.
}

% ================================================================
\begin{frame}{RQ2 --- Bigger blocks, slower generation}
  \begin{columns}[T,onlytextwidth]
    \begin{column}{0.55\textwidth}
      \panelA[\linewidth]
    \end{column}
    \begin{column}{0.43\textwidth}
      \small
      \textbf{Mean speedup, $k{=}4 \to k{=}16$}
      \begin{itemize}\itemsep0pt
        \item Qwen2.5: $2.9164 \to 1.8118$
        \item Qwen3: \; $2.6673 \to 1.0405$
      \end{itemize}
      \vspace{3pt}
      \textbf{Token timing degrades too}\\
      {\scriptsize (Qwen3, RTX4090, deterministic)}
      \begin{itemize}\itemsep0pt
        \item TTFT: \; $132.6 \to 349.6$\,ms
        \item TPOT: \; $41.2 \to 94.2$\,ms/token
      \end{itemize}
      \vspace{3pt}
      \begin{block}{}
        \small $\Beff$ \emph{rises} while latency \emph{worsens}: the
        overhead inflection arrives \alert{before $k{=}8$}.
      \end{block}
    \end{column}
  \end{columns}
\end{frame}
\note{%
\textbf{[4:55--5:55 | RQ2 the inflection]}\\[2pt]
This is the analytically interesting result, so slow down here.\\[4pt]
``Our hypothesis going in --- H2 --- was the textbook one: $k{=}8$ should beat
$k{=}4$, and $k{=}16$ might start to lose. That is wrong on this hardware. The
curve is monotonically decreasing from the smallest $k$ we tested. The
inflection has already happened before $k{=}8$.''\\[4pt]
Explain \emph{why} using the equation from earlier: the effective block size
does go up --- you really are committing more tokens per verify step --- but it
goes up sub-linearly while the drafting cost goes up linearly. Acceptance is
only 34\% at $k{=}4$ and collapses to 13\% at $k{=}16$, so most of those extra
$k$ draft steps are pure waste.\\[4pt]
The token-timing numbers on the right are the practitioner's punchline: at
$k{=}16$ you have nearly tripled time-to-first-token. For an interactive
application that is the worst possible place to spend latency --- the user sees
it directly.\\[4pt]
\emph{Implication to state:} on this class of hardware, tune $k$ downward, and
$k{=}4$ may not even be the floor.
}

% ================================================================
\begin{frame}{RQ3 --- Deterministic decoding is the safer default}
  \begin{columns}[T,onlytextwidth]
    \begin{column}{0.52\textwidth}
      \scriptsize
      \begin{tabular}{@{}lccc@{}}
        \toprule
        \textbf{Qwen3, RTX4090} & $k{=}4$ & $k{=}8$ & $k{=}16$\\
        \midrule
        Deterministic $S$ & 2.8548 & 1.9648 & 1.1091\\
        Stochastic $S$    & 2.4798 & 1.6869 & 0.9719\\
        \midrule
        Advantage         & \best{$+0.3750$} & $+0.2779$ & $+0.1372$\\
        \bottomrule
      \end{tabular}
      \vspace{8pt}
      \small
      \begin{itemize}\itemsep3pt
        \item Greedy drafting agrees with greedy verification more often
              $\Rightarrow$ higher $\alpha$ (0.336 vs.\ 0.267 at $k{=}4$).
        \item The advantage \emph{shrinks} with $k$ --- once acceptance is
              low, regime hardly matters.
      \end{itemize}
    \end{column}
    \begin{column}{0.46\textwidth}
      \begin{block}{Not universal}
        \small
        Qwen2.5 follows the same ordering at $k{=}4$ and $k{=}8$, but reverses
        slightly at $k{=}16$ ($1.8050$ det vs.\ $1.8185$ stoch).
      \end{block}
      \vspace{2pt}
      {\small\textbf{Takeaway.} If you need sampling, expect to pay
      $\sim$13\% of your speedup for it at low $k$.}
    \end{column}
  \end{columns}
\end{frame}
\note{%
\textbf{[5:55--6:30 | RQ3 regime]}\\[2pt]
Short slide --- this is the least surprising result, so move through it.\\[4pt]
The mechanism is the thing worth saying out loud: ``when both draft and target
decode greedily, the draft's proposal is compared against a deterministic
target continuation, so agreement is high. Turn sampling on for both and you
introduce two independent sources of randomness, acceptance drops from about
34\% to 27\%, and you lose roughly 13\% of your speedup.''\\[4pt]
Be careful to state the caveat rather than hiding it: at $k{=}16$ Qwen2.5
actually reverses, marginally. It is a small effect but we report it --- the
preference is a strong tendency, not a law.\\[4pt]
\emph{If running long:} this is the first slide to cut. State the one-line
takeaway --- ``deterministic wins at matched $k$, by a margin that shrinks as
$k$ grows'' --- and click through.
}

% ================================================================
\begin{frame}{RQ4 --- The same policy, a different GPU, and it all inverts}
  \begin{columns}[T,onlytextwidth]
    \begin{column}{0.53\textwidth}
      \panelB[\linewidth]
    \end{column}
    \begin{column}{0.45\textwidth}
      \small
      \textbf{Qwen3 on RTX5090-laptop}
      \scriptsize
      \begin{tabular}{@{}lcc@{}}
        \toprule
        Config & $S$ & $\alpha$\\
        \midrule
        $k{=}4$, det   & \bad{0.6743} & 0.335\\
        $k{=}4$, stoch & \bad{0.5924} & 0.285\\
        $k{=}8$, det   & \bad{0.4613} & 0.254\\
        $k{=}8$, stoch & \bad{0.3995} & 0.227\\
        $k{=}16$, det  & \bad{0.2588} & 0.180\\
        $k{=}16$, stoch& \bad{0.2222} & 0.154\\
        \bottomrule
      \end{tabular}
      \vspace{4pt}
      \small
      \begin{itemize}\itemsep1pt
        \item \alert{All six are slower than baseline.}
        \item Acceptance \emph{equals or beats} RTX4090 at every $k$ ---
              the loss is in execution cost, not the algorithm.
        \item Ordering survives: low $k$, deterministic.
      \end{itemize}
    \end{column}
  \end{columns}
\end{frame}
\note{%
\textbf{[6:30--7:50 | RQ4 portability --- the core result]}\\[2pt]
This is the slide the talk exists for. Pause before starting it.\\[4pt]
``We took the identical code, the identical models, the identical frozen
manifests, and moved to an RTX5090 laptop --- a \emph{newer} architecture. Every
single configuration is now slower than simply running the target model
autoregressively. The best case, deterministic $k{=}4$, gives us 0.67 --- we have
made inference one and a half times slower.''\\[4pt]
Then the diagnostic point, which is what makes this a finding rather than an
anecdote: ``acceptance is not the problem. At $k{=}4$ it is identical --- 0.336
against 0.335. At $k{=}8$ and $k{=}16$ the laptop actually accepts \emph{more}
--- 0.25 against 0.24, and 0.18 against 0.13. So the algorithm is doing at
least as well as on the desktop, and it is still slower everywhere. What
changed is the runtime economics: on this stack the draft model's per-step cost
relative to the target's verification cost is unfavourable, and speculative
decoding needs that ratio to be small.''\\[4pt]
Land the deployment rule: ``speedup is a property of the complete
model--hardware--runtime stack. It is not a property of the algorithm, and it is
certainly not a property of the GPU generation --- the newer card lost.''\\[4pt]
Finally, the constructive half: the \emph{ranking} did transfer. Low $k$ beats
high $k$, deterministic beats stochastic, on both devices. So policy
exploration on one machine is transferable; absolute claims are not.
}
\note{%
\textbf{[RQ4, continued]}\\[2pt]
\emph{Anticipate the question:} we do not claim to have isolated the cause. Our
candidate explanation is host-side orchestration and graph-capture behaviour ---
see the CUDA-graph note in the paper --- but we did not measure it, so we report
it as context, not as a cause. Say that before someone asks.\\[4pt]
\emph{Second likely question:} ``is the laptop card simply weaker?'' Partly ---
lower bandwidth and roughly a third of the power budget. But $S$ is a
\emph{ratio} against that same device's own autoregressive baseline, so a
uniformly slower card should not by itself push the ratio below one. The
backup hardware slide has the numbers.
}

% ================================================================
\begin{frame}{Is the effect real? Statistics and stability}
  \begin{columns}[T,onlytextwidth]
    \begin{column}{0.52\textwidth}
      \small
      \textbf{Paired one-sided $t$-tests}, $H_0{:}\ S \le 1$,\\
      paired by sample, Holm--Bonferroni across 6 configs
      (Qwen3, RTX4090, $n{=}1000$):
      \vspace{2pt}
      \scriptsize
      \begin{tabular}{@{}lcc@{}}
        \toprule
        Config & $d_S$ & $p_S$\\
        \midrule
        det $k{=}4,8,16$   & 1.18 / 1.10 / 0.29 & all $<10^{-18}$\\
        stoch $k{=}4,8$    & 1.24 / 1.13       & both $<10^{-179}$\\
        stoch $k{=}16$     & \bad{$-0.08$}     & \bad{$0.996$} \\
        \bottomrule
      \end{tabular}
      \vspace{6pt}
      \small
      \textbf{Dispersion is large.} Per-sample $S$ (det):
      $2.7561\pm0.5950$ at $k{=}4$ $\to$ $1.3549\pm0.6533$ at $k{=}16$.
      \\[2pt]
      {\scriptsize At $k{=}16$ the standard deviation is half the mean ---
      prompt-to-prompt variation dominates.}
    \end{column}
    \begin{column}{0.46\textwidth}
      \panelD[\linewidth]
      \vspace{-2pt}
      {\scriptsize Seed 999 vs.\ 123 at $k{=}4$: deterministic within
      $1.5\%$; stochastic $0.4\%$--$5.0\%$.}
    \end{column}
  \end{columns}
\end{frame}
\note{%
\textbf{[7:50--8:40 | Evidence quality]}\\[2pt]
Purpose of this slide: pre-empt the methodological objections, and show we are
not over-claiming.\\[4pt]
``Five of six RTX4090 configurations are significant after Holm--Bonferroni
correction, with large effect sizes. The sixth --- stochastic $k{=}16$ --- is
not significant, and its effect size is slightly negative, which is exactly
consistent with its sub-baseline mean. We report it as a null, not as noise.''\\[4pt]
The dispersion point matters more than the $p$-values: at $k{=}16$ the standard
deviation across prompts is roughly half the mean. A single reported average
speedup hides an enormous spread --- your slowest prompts may see no benefit at
all. This is one of the things the literature routinely omits.\\[4pt]
On the right panel: repeat seeds at $k{=}4$. Deterministic runs move by at most
1.5\% between seeds; stochastic by up to 5\%. Regime ordering is preserved
everywhere.\\[4pt]
\emph{Say the limitation before the reviewer does:} we only have repeat seeds at
$k{=}4$. We do not claim seed stability at $k{=}8$ or $k{=}16$.
}

% ================================================================
\begin{frame}{Speed is not free: quality on CNN/DailyMail}
  \centering
  \small
  \begin{tabular}{lccc}
    \toprule
    \textbf{Deterministic, $n{=}200$} & ROUGE-L & BLEU & BERTScore F1\\
    \midrule
    Autoregressive baseline & 19.58 & 3.20 & 0.793\\
    Speculative, $k{=}4$    & \bad{11.50} & \bad{1.39} & \bad{0.710}\\
    Speculative, $k{=}16$   & 14.92 & 2.04 & 0.741\\
    \bottomrule
  \end{tabular}
  \vspace{8pt}
  \begin{columns}[T]
    \begin{column}{0.47\textwidth}
      \small
      \begin{itemize}\itemsep2pt
        \item All three metrics agree: a real quality cost on the
              summarization task.
        \item It \alert{partially recovers} as $k$ grows --- moving
              \emph{against} the speed trend.
      \end{itemize}
    \end{column}
    \begin{column}{0.47\textwidth}
      \begin{block}{}
        \small The fastest configuration is also the weakest on
        summarization. There is a genuine speed--quality frontier here, not
        free acceleration.
      \end{block}
    \end{column}
  \end{columns}
\end{frame}
\note{%
\textbf{[8:40--9:25 | Quality trade-off]}\\[2pt]
Deliver this as an honest complication of our own headline number.\\[4pt]
``Earlier I said the peak was deterministic $k{=}4$. On summarization, that is
also our \emph{worst} quality point --- ROUGE-L drops from 19.6 to 11.5, and BLEU
and BERTScore move the same way. Three independent metrics agreeing means this
is not a metric artefact.''\\[4pt]
The direction is the interesting part: quality \emph{recovers} as $k$ increases,
which is the opposite direction to speed. So the two axes genuinely trade off
against each other across our grid, and the operating point is a choice, not a
default.\\[4pt]
\emph{Expected question --- be ready:} ``isn't speculative decoding supposed to
be distribution-preserving?'' Answer: exactly, and that is why this is worth
reporting. The guarantee is about the sampling distribution under the idealized
scheme; what we measure is an end-to-end implementation on a bounded generation
budget, where accepted-prefix commitment and length limits interact. We flag it
as an observed effect in our implementation and do not claim it refutes the
theory. Do not oversell this --- it is a limitation we are surfacing.\\[4pt]
Note that GSM8K and MMLU are accuracy tasks and are covered by our success
criteria; this slide is the summarization result specifically.
}

% ================================================================
\section{Wrap-up}

\begin{frame}{What we can and cannot claim}
  \begin{columns}[T,onlytextwidth]
    \begin{column}{0.48\textwidth}
      \textbf{Hypotheses}
      \small
      \begin{itemize}\itemsep2pt
        \item[\textbf{H1}] $S>1$ --- \best{supported} on RTX4090
              (11/12), \bad{rejected} on RTX5090-laptop.
        \item[\textbf{H2}] $k{=}8 > k{=}4$ --- \bad{rejected}; speedup
              falls monotonically.
        \item[\textbf{H3}] $\Beff$ up, latency worse --- \best{supported}.
        \item[\textbf{H4}] Deterministic faster --- \best{mostly supported}
              (one reversal).
      \end{itemize}
    \end{column}
    \begin{column}{0.48\textwidth}
      \textbf{Threats to validity}
      \small
      \begin{itemize}\itemsep2pt
        \item One PyTorch/Transformers stack; results do not extrapolate to
              other runtimes or GPUs.
        \item Same-family Qwen pairings only.
        \item No long-context, safety, or factual-consistency evaluation.
        \item Repeat seeds cover $k{=}4$ only.
        \item CUDA-graph capture differs across run families --- reported as
              context, \emph{not} as a measured cause.
      \end{itemize}
    \end{column}
  \end{columns}
\end{frame}
\note{%
\textbf{[9:25--9:55 | Scorecard and limits]}\\[2pt]
Move quickly --- this slide is mostly for the record and for the reviewers in
the room. Do not read every bullet.\\[4pt]
Highlight two things only. First, H2: ``we predicted $k{=}8$ would beat $k{=}4$
and it does not --- the overhead inflection is earlier than the folklore
suggests.'' Second, the limitations column exists because we want the claim
bounded: this is one runtime stack, one model family, one benchmark suite.\\[4pt]
The CUDA-graph bullet is the intellectually honest one --- we saw a difference in
capture behaviour between run families, and it is a plausible contributor to the
RTX5090 result, but we did not isolate it experimentally, so we refuse to call
it the cause.\\[4pt]
\emph{Timing checkpoint:} you should be at roughly 9:30 here. If you are past
10:30, skip straight to the takeaways.
}

% ================================================================
\begin{frame}[standout]
  \Large
  Policy ranking is portable. \\[6pt]
  Absolute speedup is not.
\end{frame}
\note{%
\textbf{[9:55--10:05 | The one line]}\\[2pt]
Pause. Let it sit for two seconds. This is the sentence you want people to
repeat to a colleague afterwards. Do not add anything to it --- click through
to the takeaways.
}

% ================================================================
\begin{frame}{Takeaways}
  \begin{enumerate}\itemsep6pt
    \item \textbf{Use a short proposal.} On both consumer GPUs and both model
          families, $k{=}4$ dominates $k{=}8$ and $k{=}16$ --- the overhead
          inflection arrives earlier than expected.
    \item \textbf{Prefer deterministic decoding} at matched $k$; sampling
          costs roughly 13\% of the speedup at $k{=}4$.
    \item \textbf{Benchmark on your target stack.} An RTX4090 gave up to
          $3.07\times$; the same code on an RTX5090 laptop was
          \emph{sub-baseline everywhere}, with unchanged acceptance.
    \item \textbf{Report dispersion and quality, not just a mean speedup.}
          Per-sample spread is large, and summarization quality moves against
          the speed trend.
  \end{enumerate}
  \vspace{6pt}
  \centering
  \small Frozen manifests, matched baselines, and raw per-sample artifacts
  support reproduction.
\end{frame}
\note{%
\textbf{[10:05--10:40 | Close]}\\[2pt]
Four sentences, one per point, then stop. Do not re-explain --- the audience has
seen all the evidence.\\[4pt]
Rule one: short proposals. Rule two: greedy if you can afford it. Rule three,
and this is the one to emphasize with your voice: measure on the machine you
will actually deploy on, because we have a concrete counter-example where a
newer GPU turned a three-times win into a one-and-a-half-times loss. Rule four:
report the spread and the quality, because the mean speedup on its own is not a
deployment decision.\\[4pt]
Close on future work in one breath: full grid across seeds, broader hardware and
model coverage, longer contexts and factuality-sensitive metrics.\\[4pt]
``Thank you --- I'm happy to take questions.''
}

% ================================================================
\appendix

\begin{frame}[standout]
  Backup slides
\end{frame}
\note{Backup material follows. Jump here only if asked.}

\begin{frame}{Backup --- full Qwen3 RTX4090 grid}
  \centering
  \small
  \begin{tabular}{lcccccc}
    \toprule
    \textbf{Config} & $S$ & $\alpha$ & $\Beff$ & $T_{\text{mean}}$ (s)
    & TTFT (ms) & TPOT (ms)\\
    \midrule
    $k{=}4$, det   & 2.8548 & 0.3362 & 1.33 & 4.8881 & 132.57 & 41.23\\
    $k{=}4$, stoch & 2.4798 & 0.2672 & 1.06 & 5.6196 & 134.54 & 45.51\\
    $k{=}8$, det   & 1.9648 & 0.2378 & 1.86 & 7.1024 & 204.99 & 58.76\\
    $k{=}8$, stoch & 1.6869 & 0.1916 & 1.50 & 8.2612 & 207.82 & 64.67\\
    $k{=}16$, det  & 1.1091 & 0.1292 & 1.96 & 12.5817 & 349.58 & 94.24\\
    $k{=}16$, stoch& 0.9719 & 0.1082 & 1.64 & 14.3390 & 354.82 & 106.03\\
    \bottomrule
  \end{tabular}
  \vspace{6pt}

  \small Target: Qwen3-4B-Instruct; draft: Qwen3-0.6B-Instruct; $n{=}1000$.
\end{frame}
\note{Use if asked for the full Qwen3 numbers, or for TTFT/TPOT detail.}

\begin{frame}{Backup --- full Qwen2.5 RTX4090 grid}
  \centering
  \small
  \begin{tabular}{lcccccc}
    \toprule
    \textbf{Config} & $S$ & $\alpha$ & $\Beff$ & $T_{\text{mean}}$ (s)
    & TTFT (ms) & TPOT (ms)\\
    \midrule
    $k{=}4$, det   & 3.0674 & 0.4212 & 1.66 & 3.7210 & 102.35 & 29.06\\
    $k{=}4$, stoch & 2.9394 & 0.4080 & 1.61 & 3.9631 & 105.26 & 30.89\\
    $k{=}8$, det   & 2.4542 & 0.3515 & 2.72 & 4.6508 & 162.76 & 36.73\\
    $k{=}8$, stoch & 2.4063 & 0.3444 & 2.67 & 4.8411 & 165.95 & 38.04\\
    $k{=}16$, det  & 1.8050 & 0.2539 & 3.76 & 6.3234 & 277.43 & 49.34\\
    $k{=}16$, stoch& 1.8185 & 0.2516 & 3.73 & 6.4059 & 278.85 & 49.81\\
    \bottomrule
  \end{tabular}
  \vspace{6pt}

  \small Target: Qwen2.5-3B-Instruct; draft: Qwen2.5-0.5B-Instruct; $n{=}1000$.\\
  Higher $\alpha$ and $\Beff$ than Qwen3 at every $k$.
\end{frame}
\note{%
Qwen2.5 accepts more than Qwen3 at every $k$ --- $\alpha=0.42$ vs.\ $0.34$ at
$k{=}4$, and effective block size 3.76 vs.\ 1.96 at $k{=}16$. That is why
Qwen2.5 never falls below baseline on the RTX4090 while Qwen3 does. Likely a
smaller capability gap between the 3B target and 0.5B draft than between the
4B target and 0.6B draft.
}

\begin{frame}{Backup --- acceptance portability}
  \centering
  \panelC[0.62\linewidth]
  \vspace{2pt}

  \small Qwen3 acceptance on the laptop is identical at $k{=}4$ (0.335 vs.\
  0.336) and \emph{higher} at $k{=}8$ and $k{=}16$ (0.254 vs.\ 0.238; 0.180
  vs.\ 0.129). The algorithm behaves at least as well; the
  \alert{runtime economics} differ.
\end{frame}
\note{%
The single most useful backup slide for the RQ4 discussion. If someone
challenges the RTX5090 result --- ``did you misconfigure it?'' --- this is the
answer: acceptance is a model-level quantity; it is unchanged at $k{=}4$ and
actually \emph{better} at higher $k$, so the speculative loop is doing at least
as well as it did on the desktop. The loss is in per-step execution cost, not
in the algorithm.
}

\begin{frame}{Backup --- success criteria and metrics}
  \begin{columns}[T,onlytextwidth]
    \begin{column}{0.48\textwidth}
      \small
      \textbf{Stated success criteria}
      \begin{itemize}\itemsep2pt
        \item Primary: mean $S \ge 1.3\times$ with Holm--Bonferroni-corrected
              $p_S < 0.05$.
        \item Primary: absolute quality drop $\le 1.0$ point on GSM8K and
              MMLU, deterministic regime.
        \item Secondary: best $(k, \text{regime})$ under quality constraints;
              dispersion; seed stability; cross-hardware ranking transfer.
      \end{itemize}
    \end{column}
    \begin{column}{0.48\textwidth}
      \small
      \textbf{Key metric definitions}
      \begin{itemize}\itemsep2pt
        \item $S$ = baseline latency / speculative latency
        \item $\alpha$ = accepted / proposed draft tokens
        \item $\Beff$ = accepted draft tokens / verify steps
        \item TTFT = first token time $-$ request start
        \item TPOT = (completion $-$ first token) / output tokens
        \item $d_S$ = Cohen's $d$, paired baseline vs.\ speculative latency
      \end{itemize}
    \end{column}
  \end{columns}
\end{frame}
\note{Reference slide for methodology questions.}

\begin{frame}{Backup --- hardware profiles}
  \centering
  \small
  \begin{tabular}{lll}
    \toprule
    \textbf{Component} & \textbf{RTX4090 desktop} & \textbf{RTX5090 laptop}\\
    \midrule
    GPU memory        & 24\,GB GDDR6X & 24\,GB GDDR7\\
    Memory bandwidth  & \best{1008\,GB/s} & 896\,GB/s\\
    CUDA cores        & \best{16384} & 10496\\
    Architecture      & Ada Lovelace & \best{Blackwell}\\
    Power profile     & \best{450\,W TGP class} & up to 150\,W TGP (OEM-dependent)\\
    Software stack    & \multicolumn{2}{c}{PyTorch + Transformers, pinned per run family}\\
    \bottomrule
  \end{tabular}
  \vspace{8pt}

  \small The laptop part is the newer architecture but has lower bandwidth and
  roughly \alert{one third of the power budget}.
\end{frame}
\note{%
The obvious question after RQ4 is ``is the 5090 laptop just a weaker card?''
Partly --- lower bandwidth, far lower sustained power. But note that this affects
the autoregressive baseline too, and $S$ is a \emph{ratio} against that same
device's baseline. So a uniformly slower card should not by itself flip the
ratio below 1. That is why we point at the draft/target cost ratio and the
runtime path rather than at raw card performance.
}

\end{document}
